% concatenated from main.tex
\documentclass{article}
\usepackage[top=1.0in, left=1.2in, right=1.2in, bottom=1.0in]{geometry}
\usepackage[numbers,sort&compress]{natbib}
\usepackage{graphicx}
\usepackage{amsmath}
\usepackage{dsfont}
\usepackage{amssymb}
\usepackage{setspace}
\usepackage{textgreek}
\usepackage{titlesec}
\usepackage{caption}
\usepackage{xcolor}
\usepackage[most]{tcolorbox}
\usepackage[colorlinks=true,
            linkcolor=blue,
            citecolor=blue,
            urlcolor=blue,
            pdfborder={0 0 0}]{hyperref}
\usepackage[version=4,arrows=pgf-filled,textfontname=sffamily,mathfontname=mathsf]{mhchem}

\usepackage{setspace}
\onehalfspacing

\usepackage{amsthm}
\theoremstyle{plain}
\newtheorem{theorem}{Theorem}[section]
\newtheorem{proposition}{Proposition}[section]
\newtheorem{corollary}{Corollary}[section]
\newtheorem{lemma}{Lemma}[section]
\theoremstyle{definition}
\newtheorem{remark}{Remark}[section]
\newtheorem{assumption}{Assumption}[section]
\newtheorem{definition}{Definition}[section]
\usepackage{authblk}
\renewcommand{\vec}[1]{\boldsymbol{#1}}
\newcommand{\rev}[1]{{\color{red}#1}}
\hyphenation{observations consider}

\title{Structural functional identifiability and model discovery in \\ differential equation models}
\author[1]{Torkel E Loman}
\author[2]{Alexander P Browning}
\author[1]{Ruth E Baker}
\date{}

\affil[1]{Mathematical Institute, University of Oxford, Oxford, United Kingdom}
\affil[2]{School of Mathematics and Statistics, University of Melbourne, Australia}


% Build setup: keep a single source file in the workspace.

\begin{document}

\maketitle

% \tableofcontents

\begin{abstract}

Differential equation models are widely used to describe, interpret, and predict dynamical phenomena across science and engineering. In practice, however, the governing dynamics are rarely fully known and must be inferred from observational data. Traditionally, inverse problems in differential equation modelling have focused on estimating unknown parameter values. In this setting, structural identifiability determines whether parameter values can, in principle, be uniquely recovered from ideal observations and is, therefore, a prerequisite for meaningful inference. More recently, the integration of machine learning with mechanistic modelling has enabled the discovery of unknown equations, functions, and constitutive relationships, substantially expanding the space of admissible models. This raises a fundamental question: under what conditions can unknown functional components be uniquely recovered from data? In this paper, we generalise the classical notion of structural parameter identifiability to functional identifiability. We first identify broad classes of models for which unique functional recovery is impossible. We then show how functional identifiability can be assessed for differential equation models using differential algebra-based techniques which are well-established as a means of assessing structural identifiability for ordinary differential equation-based models. Our framework reveals new phenomena that arise in the transition from parametric to functional inference and have no analogue in the classical setting. Finally, we characterise functional identifiability in several common model classes. Taken together, our results demonstrate that functional identifiability provides a theoretical foundation for modern inverse problems in differential equation modelling, particularly those that use machine learning representations of unknown system components.

\end{abstract}

% ------------------------------------------------------------------------------------------------------------------------------------ %
% ----------------------------------------------------------- Introduction ----------------------------------------------------------- %
% ------------------------------------------------------------------------------------------------------------------------------------ %

\section{Introduction}

Mathematical models are widely used to describe, explain, and predict the behaviour of dynamical systems. Traditionally, such models are mechanistic: their governing equations are derived from physical, biological, or other domain-specific principles. In practice, however, these principles rarely determine a model completely. Mechanistic models, therefore, typically contain unknown quantities that must be inferred from observational data before the models can be used for prediction or decision-making. Historically, these unknowns have been finite-dimensional parameters, leading to a rich literature on parameter estimation and identifiability~\cite{bellman_structural_1970,raue_structural_2009,Chis_struct_ident_2011,villaverde_protocol_2022,simpson_parameter_2026}. Recent advances in machine learning have substantially expanded the range of unknown model components that can be inferred from data. Rather than estimating only scalar parameters, it is now increasingly common to learn unknown functional relationships directly from observations. This has led to the emergence of hybrid modelling approaches that combine mechanistic structures with learning of unknown functions and have found applications across the physical, biological, and engineering sciences~\cite{alber_integrating_2019,noordijk_rise_2024,willard_integrating_2022}.

In this work, we focus on \emph{function inference}, where the structure of a differential equation model is assumed known up to one or more unknown functions that must be learnt from data (Figure~\ref{fig:ude_workflow}). These functions may represent unknown rate laws, interaction terms, transfer processes, feedback relationships, potential fields, or other system components with forms that are not known \textit{a priori}. Unlike scalar parameters, functions are infinite-dimensional objects, substantially expanding the space of admissible models. In practice, such functions are often represented using flexible approximators such as neural networks or Gaussian processes~\cite{rackauckas_universal_2021,raissi_physics-informed_2019,karniadakis_physics-informed_2021,dandekar_bayesian_2022}. This shift from parameter inference to function inference raises a natural question: under what conditions can an unknown function be uniquely recovered from data?

\begin{figure}
    \centering
    \includegraphics[width=0.95\linewidth]{Figures/ude_workflow.pdf}
    \caption{\textbf{Typical function learned workflow.} Example workflow where a functional form is learnt using a universal differential equation combined with data. (A) A differential equation model is formulated, including unknown parameters (here $c$ and $d$) and functions (here $f$). (B) The unknown function is replaced with a universal function approximator (here, a neural network) which represents an arbitrary function through a finite set of parameters. (C) The model's scalar parameters (both mechanistic parameters and those representing the unknown function) are fitted to data using a standard parameter fitting workflow. (D) The fitted model can be evaluated to, for example, predict the time evolution of unobserved variables (here, $X_2(t)$) or the fitted function (by evaluating the fitted function approximator on some domain).}
    \label{fig:ude_workflow}
\end{figure}

In the parameter estimation literature, \emph{identifiability} concerns whether model parameters can be uniquely determined from observational data~\cite{bellman_structural_1970,cobelli_parameter_1980,hines_determination_2014,villaverde_protocol_2022}. A distinction is commonly drawn between \emph{structural} and \emph{practical} identifiability. Methods for structural identifiability analysis ask whether distinct parameter values can produce identical observable dynamics and thus remain indistinguishable even with unlimited noiseless data~\cite{bellman_structural_1970,chis_structural_2011,meshkat_algorithm_2009}. Practical identifiability, by contrast, concerns whether parameters can be reliably estimated from the finite and noisy data available in practice~\cite{raue_structural_2009,simpson_parameter_2026,wieland_structural_2021}. Structural identifiability provides the theoretical foundation for parameter inference, and hence it should be assessed before both a practical identifiability analysis and model inference.

The shift from parameter inference to functional inference raises a natural analogue of this classical problem. We recently introduced \emph{functional identifiability} as a framework for determining whether unknown functions can be uniquely recovered from data~\cite{loman_functional_2025}. Here we develop a theory of \emph{structural functional identifiability} for ordinary differential equation (ODE) models containing unknown functions. Our contributions are fourfold. First, we identify broad classes of models in which functional non-identifiability arises naturally and show that, in these cases, hybrid models are theoretically equivalent to neural differential equations with entirely unknown right-hand sides~\cite{chen_neural_2019,kidger_neural_2020}. Second, we extend differential algebra-based methods to assess functional identifiability. Third, we demonstrate that structural functional identifiability exhibits phenomena with no direct analogue in the parametric setting, including entanglement between unknown functions and parameters, entanglement between multiple unknown functions, and potential non-identifiability where the single unknown component is a function. Finally, we illustrate these methods on commonly used model types, including a chemical reaction network model and the Lotka--Volterra model.

% ------------------------------------------------------------------------------------------------------------------------------------ %
% ----------------------------------------------------------- SI ---------------------------------------------------------------- %
% ------------------------------------------------------------------------------------------------------------------------------------ %

\section{Structural functional identifiability: definition and scope of this work} \label{sec:functional_identifiability_definition}

Structural functional identifiability asks whether an unknown function can, in theory, be determined from perfect observations of a model's output. Consider an ODE model
\begin{align*}
    \dot{\vec{X}}(t)&= \vec{F}(\vec{X};\, \vec{f},\, \vec{p}), \\
    \vec{Y}(t)&=\vec{H}(\vec{X}(t)),
\end{align*}
where $\vec{X}\in\mathbb{R}^n$ is the state vector, $\vec{Y}$ is the observed output, $\vec{F}$ and $\vec{H}$ are known functions, $\vec{p}\in\mathbb{R}^m$ is a vector of unknown scalar parameters, and $\vec{f}=(f_1,\ldots,f_k)$ is a collection of unknown functions with
$f_j:\mathcal{U}_j\to\mathbb{R}$  and $\mathcal{U}_j\subseteq\vec{X}$. Two choices of the unknowns, $(\vec{f}_1,\vec{p}_1)$ and $(\vec{f}_2,\vec{p}_2)$, are said to give the same observed outputs if they produce the same output trajectory $\vec{Y}(t)$ for every initial condition consistent with prior knowledge of the initial conditions. Often, the known initial conditions are simply $\vec{Y}(0)$, however, additional initial condition knowledge is possible, and is briefly considered in the example in Section~\ref{sec:diffalg_two_parameter}. Throughout this work, we will assume that $\vec{Y}$ consists of some or all of the state variables in $\vec{X}$. More complicated observable formulae, however, would also be possible to consider.

\begin{definition} \label{def:functional_identifiability}
Fix a generic choice of the unknowns $(\vec{f}_1,\vec{p}_1)$. An unknown function $f_j$ is:
\begin{itemize}
    \item \emph{Globally structurally identifiable} if every choice
    of unknowns $(\vec{f}_2,\vec{p}_2)$ that gives the same observed outputs as
    $(\vec{f}_1,\vec{p}_1)$ has $f_{1,j}=f_{2,j}$.
    \item \emph{Locally structurally identifiable} if there exists
    $\varepsilon>0$ such that every choice of unknowns $(\vec{f}_2,\vec{p}_2)$ that gives the
    same observed outputs as $(\vec{f}_1,\vec{p}_1)$ and satisfies
    $\sup_{\vec{u}\in\mathcal{U}_j}|f_{1,j}(\vec{u})-f_{2,j}(\vec{u})|<\varepsilon$
    must have $f_{1,j}=f_{2,j}$. Often, but not always, this is equivalent to there only being a finite number of functions $f_j$ giving the same observed dynamics.
    \item \emph{Structurally non-identifiable} otherwise.
\end{itemize}
\end{definition}

\noindent
Throughout this work we will mainly consider global structural identifiability, for simplicity termed structural identifiability unless otherwise mentioned. Local structural identifiability will be considered in more detail in Section~\ref{sec:local_global}.

Determining whether a model's components are structurally identifiable has both practical and theoretical applications. Primarily, it is used to determine what model components can be recovered from data. However, it can also guide experimental design: by investigating how structural identifiability depends on the observed quantities, one can better determine what system components to measure to maximise identifiability.

This paper is laid out as follows. Section~\ref{sec:general_nonident} identifies several broad classes of models for which structural functional non-identifiability holds generally. The aim is not to provide a comprehensive catalogue of non-identifiable models, but to highlight common cases where functional recovery is an ill-posed problem. Section~\ref{sec:diffalg} then illustrates, through a few examples, how differential algebra techniques can be used to determine functional identifiability in ODE models. This approach applies beyond the special forms considered in Section~\ref{sec:general_nonident}, and can, therefore, handle models with diverse algebraic structures. Section~\ref{sec:nonidentifiability_types} uses the same approach to distinguish different forms of functional non-identifiability that can be encountered, considering both global and local structural identifiability. Finally, Section~\ref{sec:applications} applies the differential algebra approach to two common model categories, including cases where unknown functions appear multiple times in the same model or depend on multiple variables.

% ------------------------------------------------------------------------------------------------------------------------------------ %
% ------------------------------------- Results 1: General functional non-identifiability results ------------------------------------ %
% ------------------------------------------------------------------------------------------------------------------------------------ %

\section{Structural non-identifiability for general classes of models} \label{sec:general_nonident}

In this section, we identify broad classes of ODE models for which unique functional recovery is structurally impossible. For each class, we construct an explicit family of functions that produces identical observable dynamics and is, therefore, functionally indistinguishable. Throughout, we assume that all state variables are observable, so that $\vec{Y}(t)=\vec{X}(t)$. Consequently, the same non-identifiability results also apply when only a subset of variables is observed. Proofs of each result are provided in Appendix~\ref{appendix:proofs_section2}. We note that the equational forms considered here are by no means exhaustive (for example, here unknown parameters and functions occur in separate terms, which is not required). In the next section, we will show how to determine structural identifiability more generally.

% --------- Single-variable ODEs --------- %

\subsection{Scalar ODE models} \label{sec:nonident_1var_UDEs}

Consider a scalar ODE model containing both known and unknown dynamics:
\begin{equation} \label{eq:nonident_1var_UDEs}
    \dot{X} = f(X) + g(X;\vec{p}),
\end{equation}
where $f:\mathbb{R}\to\mathbb{R}$ is an unknown function and $g:\mathbb{R}\to\mathbb{R}$ is a known functional form depending on a vector $\vec{p}$ with $|\vec{p}|>0$ (i.e. at least one unknown parameter is present).

\begin{proposition} \label{prop:1var_nonident}
The unknown function $f$ is structurally non-identifiable. Specifically, given any admissible pair $(f_1, \vec{p}_1)$, any pair $(f_2, \vec{p}_2)$ of the form
\begin{align}    
    \vec{p}_2 &\in \mathbb{R}^{|\vec{p}|}, \\
    f_2(X) &= f_1(X) + g(X;\vec{p}_1) - g(X;\vec{p}_2), \label{eq:1var_nonident_expr}
\end{align}
yields identical dynamics to $(f_1, \vec{p}_1)$. 
\end{proposition}

The non-identifiability arises because the unknown function $f$ can absorb arbitrary changes in the parametric term $g(X;\vec{p})$. Observable data therefore constrain only the combined dynamics $f(X)+g(X;\vec{p})$, rather than the individual contributions of the known and unknown components. Consequently, the mechanistic component $g$ and the learned component $f$ cannot be disentangled from observations alone: any pair $(f_2,\vec{p}_2)$ belonging to the family described by Proposition~\ref{prop:1var_nonident} yields identical dynamics, making simultaneous recovery of the parameter values and the unknown function structurally impossible. Example models satisfying the conditions of Proposition~\ref{prop:1var_nonident} are provided in Appendix~\ref{appendix:examples_1var}, with Figure~\ref{fig:single_var_nonident} providing a simple example.

\begin{figure}
    \centering
    \includegraphics[width=0.90\linewidth]{Figures/self_activation_loop_hill_functions.pdf}
    \caption{\textbf{A simple self-activation loop model is inherently structurally non-identifiable.} (A) A self-activation loop where a single component, $X$, activates its own production according to an unknown function $f(X)$ and decays at a linear (unknown) rate $d$. According to Proposition~\ref{prop:1var_nonident}, this model is structurally non-identifiable. The reason is that, given any base pair $(f_1(X),d_1)$, then any alternative pair $d_2 \in \mathbb{R}$, $f_2(X) = f_1(X) + (d_2 - d_1)X$ will yield identical dynamics. (B) Using a sigmoidal base function $f_1(X)$ (black dashed line), we plot potential functions $f_2(X)$ for different choices of $d_2$(coloured lines). (C) For each pair $(f_2(x),d_2)$ in B, we simulate the ODE in A (coloured lines) and compare these to the simulation for $(f_1(x),d_1)$. These simulations confirm that the different model alternatives all yield identical dynamics. While we here only consider a single initial condition, Proposition~\ref{prop:1var_nonident} tells us that this will hold for any initial condition.}
    \label{fig:single_var_nonident}
\end{figure}


% --------- Multi-variable ODEs --------- %

\subsection{Partially augmented systems of ODEs} \label{sec:nonident_partial_system_UDEs}

The previous example can be extended to systems of ODEs by adding additional variables not directly involved in the non-identifiability relation. As a concrete example, consider an oscillator with position $X$ and velocity $\dot{X}$,
\begin{align*}
    \dot{X} &= X, \\
    \ddot{X} &= f(X) + kX - c\dot{X},
\end{align*}
where $f:\mathbb{R}\to\mathbb{R}$ is an unknown restoring-force contribution and $k>0$ and $c>0$ are unknown stiffness and damping parameters. The acceleration depends on the sum $f(X)+kX$, where changes $k$ can be offset by tuning a linear term in $f(X)$.

We can generalise this model structure by considering the following system in which one equation contains an unknown function:
\begin{align}
    \dot{X}_1 &= f(\vec{u}) + g_1(\vec{X};\vec{p}_1),  \label{eq:nonident_partial_system_UDEs} \\
    \dot{X}_i &= g_i(\vec{X};\vec{p}_i), \qquad i = 2,\ldots,n, \notag
\end{align}
{where $\vec{u} \subseteq \vec{X}$, $f$ is unknown, the $g_i$ are known functional forms depending on unknown parameter sets $\vec{p}_i$, and $|\vec{p}_1| > 0$. Suppose further that $g_1$ can be decomposed as
\begin{equation*}
    g_1(\vec{X};\vec{p}_1) = r(\vec{v};\vec{a}) + s(\vec{X};\vec{b}),
\end{equation*}
where $\vec{v} \subseteq \vec{u}$, $\vec{a} \cap \vec{b} = \emptyset$, $\vec{a} \cap \vec{p}_i = \emptyset$ for all $i > 1$, and $|\vec{a}| > 0$. In practice, this means that the known part of the first equation contains a parameter-dependent term acting on variables that are arguments of the unknown function $f$. Furthermore, the parameters of this term cannot occur in any other equation.

\begin{proposition} \label{prop:multivar_nonident}
Under these conditions, $f$ is structurally non-identifiable. Specifically, given any admissible pair $(f_1, \vec{p}_{1,1}) \equiv (f_1, \{\vec{a}_{1}, \vec{b}\})$, any pair $(f_2, \vec{p}_{1,2}) \equiv (f_2, \{\vec{a}_{2}, \vec{b}\})$ of the form
\begin{align}
    \vec{a}_2 &\in \mathbb{R}^{|\vec{a}|}, \notag \\
    \label{eq:multivar_nonident_expr}
    f_2(\vec{u}) &= f_1(\vec{u}) + r(\vec{v};\vec{a}_1) - r(\vec{v};\vec{a}_2),
\end{align}
yields identical dynamics to $(f_1, \vec{p}_{1,1})$.
\end{proposition}

As in Proposition~\ref{prop:1var_nonident}, the non-identifiability arises because changes in the known component of the equation for $X_1$ can be absorbed into the unknown function without altering the observable dynamics. Example models satisfying the conditions of Proposition~\ref{prop:multivar_nonident} are provided in Appendix~\ref{appendix:examples_multivar}. We note that this result is only guaranteed to hold when $f$ appears only in the equation for $X_1$; occurrences of $f$ in the equations for $X_2,\ldots,X_n$ can introduce additional constraints and may restore identifiability, as shown in Appendix~\ref{appendix:repeated_function_counterexample}. 

% --------- Complete UDEs --------- %

\subsection{Fully augmented systems of ODEs} \label{sec:nonident_complete_UDEs}

Finally, we consider the case of a system of fully augmented ODEs where every equation contains an additive unknown function of the full state vector:
\begin{equation}
    \label{eq:complete_UDE}
    \dot{X}_i = f_i(\vec{X}) + g_i(\vec{X};\vec{p}_i), \qquad i = 1,\ldots,n,
\end{equation}
where the $f_i$ are equation-specific, and $|\vec{p}_i|>0$, that is, every equation contains at least one unknown parameter. Note that the different $\vec{p}_i$ need not be disjoint, and may overlap partially or completely.

\begin{proposition} \label{prop:complete_UDE_nonident}
All unknown functions $f_i$ are structurally non-identifiable. Specifically, given any admissible pair $(f_{i,1}, \vec{p}_{i,1})$, any pair of the form
\begin{align}
    \vec{p}_{i,2} &\in \mathbb{R}^{|\vec{p}_i|}, \notag \\
    \label{eq:complete_UDE_nonident_expr}
    f_{i,2}(\vec{X}) &= f_{i,1}(\vec{X}) + g_i(\vec{X};\vec{p}_{i,1}) - g_i(\vec{X};\vec{p}_{i,2}),
\end{align}
with the choices of $\vec{p}_{i,2}$ required to be consistent on any overlaps between the $\vec{p}_i$, yields dynamics identical to $(f_{i,1}, \vec{p}_{i,1})$ for all $i$.
\end{proposition}

For example, a microbial biomass $X$ growing on a nutrient $S$ could be written
\begin{align*}
    \dot{X} &= f_X(X,S) - dX - cX^2, \\
    \dot{S} &= -f_S(X,S) - \lambda S,
\end{align*}
where $f_X:\mathbb{R}^2\to\mathbb{R}$ and $f_S:\mathbb{R}^2\to\mathbb{R}$ are unknown growth and nutrient-consumption laws, while $d$, $c$, and $\lambda$ are unknown decay parameters. Even with both $X$ and $S$ observed, $d$ and $c$ can be tuned arbitrarily with $f_X$ being adjusted accordingly to preserve the dynamics of $X$, and similarly for $f_S$ and $\lambda$. Further examples are provided in Appendix~\ref{appendix:examples_complete}. 

% --------- Neural ODEs --------- %

\subsection{Neural ODEs} \label{sec:nonident_neural_odes}

Finally, we consider a system of ODEs where each right-hand side is an unknown function of the full state space, i.e. no mechanistic knowledge is provided. These models are also known as \emph{neural ODEs}~\cite{chen_neural_2019,kidger_neural_2020}, and can be written as 
\begin{equation} \label{eq:nonident_neural_odes}
    \dot{X}_i = f_i(\vec{X}),
\end{equation}
where each $f_i$ is an unknown function.

\begin{proposition} \label{prop:nonident_neural_odes}
Neural ODEs are structurally identifiable if and only if the full state vector, $\vec{X}$, is observable.
\end{proposition}

This result suggests that, for fully observable ODE systems, the full right-hand side is always structurally identifiable. This is particularly notable, as a neural ODE is the baseline, fully data-driven, case where no mechanistic knowledge is provided. Here, in several non-identifiable models (including those described in Section~\ref{sec:nonident_1var_UDEs} and Section~\ref{sec:nonident_complete_UDEs}), the identifiable expression(s) are exactly the set of ODE right-hand side expressions. This implies that, for these model cases, when a hybrid model is proposed, no actual mechanistic knowledge is provided. Here, fitting the ODE model incorporating unknown functions to data is equivalent to fitting a fully data-driven neural ODE. This highlights a practical role for structural functional identifiability: determining when mechanistic assumptions genuinely provide identifiable information and when they do not.

The results in this section identify broad settings where functional recovery fails by construction. Next, we will describe how the differential algebra approach provides a practical framework for determining whether generic, partially observed, models are structurally identifiable.

% ------------------------------------------------------------------------------------------------------------------------------------ %
% ---------------------------- Results 2: Determining structural identifiability with differential algebra --------------------------- %
% ------------------------------------------------------------------------------------------------------------------------------------ %

\section{Determining structural identifiability via the differential algebra approach} \label{sec:diffalg}

% \subsection{Differential algebra and input-output equations} \label{sec:diffalg_intro}

Differential algebra provides a way to remove unobserved state variables from an ODE model through repeated differentiation and substitution. The result is an \emph{input-output equation}: a differential equation involving only observed variables, their derivatives, and the unknown functions and parameters. If two choices of the unknowns generate the same input-output equation, then they cannot be distinguished given the observed dynamics (implying structural non-identifiability). Crucially, structural identifiability considers the case of arbitrary amounts of noiseless data, which implies that each observable and each of their differentials can be observed independently. In practice, this means that one can consider each term of the input-output equation in isolation, from which identifiable quantities can be recovered.

Consider, for example, the following system of ODEs,
\begin{equation*}
    \dot X_1=-aX_1+bX_2,\qquad
    \dot X_2=cX_1-dX_2,
\end{equation*}
with only $X_1$ observed and where $a$, $b$, $c$, and $d$ are unknown parameters. Eliminating the unobserved variable $X_2$ by differentiating and substituting gives the input-output equation
\begin{equation*}
    \ddot X_1+(a+d)\dot X_1+(ad-bc)X_1=0.
\end{equation*}
Here, $\dot{X_1}$ can, in the theoretical case of ideal data, be varied independently of $\ddot{X_1}$ and $X_1$. From this, the sum $(a+d)$ can be identified. Similarly, $ad-bc$ is identifiable. However, as multiple choices of $a$, $b$, $c$, and $d$ preserve these coefficients, the individual parameters are not structurally identifiable from observations of $X_1$ alone. If instead $d$ is known, $a$ is identifiable (as well as the product $bc$, but not the individual parameters $b$ and $c$).

% --------- One-parameter mutual activation model --------- %

\subsection{Single-parameter mutual activation loop model} \label{sec:diffalg_one_parameter}

To demonstrate how differential algebra can be applied to determine functional identifiability, we first consider a mutual activation loop model. Here, two variables mutually activate each other and decay at the same constant rate:
\begin{align}
    \dot{X}_1 &= f(X_2) - dX_1, \label{eq:diffalg_model1_dXdt} \\
    \dot{X}_2 &= X_1 - dX_2,   \label{eq:diffalg_model1_dYdt}
\end{align}
with $f:\mathbb{R}\to\mathbb{R}$ an unknown function and $d$ an unknown parameter.

\begin{proposition} \label{prop:diffalg_one_parameter_identifiability}\label{prop:diffalg_case1}\label{prop:diffalg_case2}
In the single-parameter mutual activation model, Equations~\eqref{eq:diffalg_model1_dXdt}--\eqref{eq:diffalg_model1_dYdt}, both $f$ and $d$ are structurally identifiable in either of the following cases:
\begin{enumerate}
    \item Only $X_2$ is observable.
    \item Only $X_1$ is observable, provided $f$ is locally invertible and $d \neq 0$.
\end{enumerate}
\end{proposition}

\begin{proof}[Proof of part (i)]
To eliminate $X_1$, we first isolate it in Equation~\eqref{eq:diffalg_model1_dYdt} and differentiate:
\begin{equation*}
    X_1= \dot{X}_2 + dX_2 \quad \Longrightarrow \quad \dot{X}_1 = \ddot{X}_2 + d\dot{X}_2.
\end{equation*}
Substituting into Equation~\eqref{eq:diffalg_model1_dXdt} gives the following input-output equation
\begin{equation}
    \label{eq:diffalg_model1_Yobs_dY2dt2}
    \ddot{X}_2 + 2d\dot{X}_2 + d^2 X_2 - f(X_2) = 0.
\end{equation}
Now suppose two pairs $(f_1, d_1)$ and $(f_2, d_2)$ yield identical $X_2$ dynamics. Inserting both into Equation~\eqref{eq:diffalg_model1_Yobs_dY2dt2}, isolating $\ddot{X_2}$, and then equating gives
\begin{equation}
    \label{eq:diffalg_model1_Yobs_equiveq}
    -2d_1\dot{X}_2 - d_1^2 X_2 + f_1(X_2) = \ddot{X}_2 = -2d_2\dot{X}_2 - d_2^2 X_2 + f_2(X_2).
\end{equation}
Here, under the assumption of an arbitrary amount of noiseless data, each of $X_2$'s derivatives can be varied independently, allowing us to separate Equation~\eqref{eq:diffalg_model1_Yobs_equiveq} into two independent equalities that must hold:
\begin{align}
    -2d_1\dot{X}_2       &= -2d_2\dot{X}_2,       \label{eq:diffalg_model1_Yobs_split1} \\
    -d_1^2 X_2 + f_1(X_2) &= -d_2^2 X_2 + f_2(X_2).  \label{eq:diffalg_model1_Yobs_split2}
\end{align}
Equation~\eqref{eq:diffalg_model1_Yobs_split1} immediately gives $d_1 = d_2 \equiv d$. Substituting into Equation~\eqref{eq:diffalg_model1_Yobs_split2} then gives $f_1(X_2) = f_2(X_2)$. This demonstrates that two different pairs $(f, d)$ cannot yield identical dynamics with respect to $X_2$, showing that both $d$ and $f$ are structurally identifiable.
\end{proof}

The proof of part~(ii) is analogous to the proof above, but with $X_2$ eliminated rather than $X_1$. Since $f$ has the unobserved variable $X_2$ as its argument, eliminating $X_2$ requires inverting $f$. Crucially, only local invertibility is needed; this means that identifiability can hold for functions such as $f(X_2)=\sin(X_2)$ that are not globally invertible. The proof of part~(ii) is given in Appendix~\ref{appendix:proof_case2}.

% --------- Two-parameter mutual activation model --------- %

\subsection{Two-parameter mutual activation loop model} \label{sec:diffalg_two_parameter}

We next allow the two variables to have different decay rates, giving
\begin{align}
    \dot{X}_1 &= f(X_2) - cX_1, \label{eq:diffalg_model2_dXdt} \\
    \dot{X}_2 &= X_1 - dX_2,   \label{eq:diffalg_model2_dYdt}
\end{align}
where $f:\mathbb{R}\to\mathbb{R}$ is an unknown function and $c$ and $d$ are unknown parameters.

\begin{proposition} \label{prop:diffalg_two_parameter_identifiability}\label{prop:diffalg_case3}\label{prop:diffalg_case4}
In the two-parameter mutual activation model, Equations~\eqref{eq:diffalg_model2_dXdt}--\eqref{eq:diffalg_model2_dYdt}, the identifiability properties are as follows.
\begin{enumerate}
    \item If only $X_1$ is observable then, provided $f$ is locally invertible and $c \neq 0$: if $f^{-1}$ is non-affine, then $f$, $c$, and $d$ are globally structurally identifiable. If $f^{-1}$ is affine, then $c$, $d$, and $f$ are each only locally identifiable (the slope of $f^{-1}$, however, is globally identifiable).
    \item If only $X_2$ is observable then $f$ is structurally non-identifiable and only the quantities $c + d$ and $f(X_2) - cdX_2$ are identifiable. Specifically, given any admissible triplet $(f_1, c_1, d_1)$, any alternative choice of the form
\begin{align}
    \label{eq:diffalg_case4_nonident_expr}
    c_2 &\in\mathbb{R}, \notag  \\
    d_2    &= c_1 + d_1 - c_2, \notag  \\
    f_2(X_2) &= f_1(X_2) + (c_2 c_1 + c_2 d_1 - c_2^2 - c_1 d_1)X_2,
\end{align}
    yields dynamics identical to $(f_1, c_1, d_1)$ with respect to $X_2$. Here, specifically, the identifiable quantities are $c+d$ and $f(x) - cdx$
\end{enumerate}
\end{proposition}

In part~(i), when $f^{-1}$ is affine, the values of $c$ and $d$ can be recovered but cannot be assigned to the two decay processes: the unordered pair $\{c,d\}$ is identifiable, while the permutation is not. The proof of part~(i) is given in Appendix~\ref{appendix:proof_case3}. Global and local identifiability are discussed further in Section~\ref{sec:local_global}.

\begin{proof}[Proof of part (ii)]
To eliminate $X_1$, we first isolate it in Equation~\eqref{eq:diffalg_model2_dYdt} and differentiate:
\begin{equation*}
    X_1= \dot{X}_2 + dX_2 \quad \Longrightarrow \quad \dot{X}_1 = \ddot{X}_2 + d\dot{X}_2.
\end{equation*}
Substituting into Equation~\eqref{eq:diffalg_model2_dXdt} gives the following input-output equation
\begin{equation}
    \label{eq:diffalg_model2_Yobs_Yonly}
    \ddot{X}_2 = -(c+d)\dot{X}_2 + f(X_2) - cdX_2.
\end{equation}
Now suppose two pairs $(f_1, d_1)$ and $(f_2, d_2)$ yield identical $X_2$ dynamics. Inserting both into Equation~\eqref{eq:diffalg_model2_Yobs_Yonly}, isolating $\ddot{X_2}$, and then equating gives
\begin{equation}
    \label{eq:diffalg_model2_Yobs_equiveq}
    -(c_1+d_1)\dot{X}_2 + f_1(X_2) - c_1 d_1 X_2 = \ddot{X}_2 = -(c_2+d_2)\dot{X}_2 + f_2(X_2) - c_2 d_2 X_2.
\end{equation}
Here, each of the derivatives of $X_2$ can be varied independently, allowing us to separate Equation~\eqref{eq:diffalg_model2_Yobs_equiveq} into two independent equalities:
\begin{align}
    c_1 + d_1 &= c_2 + d_2, \label{eq:diffalg_case4_sum} \\
    f_1(X_2) - c_1 d_1 X_2 &= f_2(X_2) - c_2 d_2 X_2. \label{eq:diffalg_case4_fY}
\end{align}
For any $c_2 \in \mathbb{R}$, setting $d_2 = c_1 + d_1 - c_2$ satisfies Equation~\eqref{eq:diffalg_case4_sum}. Equation~\eqref{eq:diffalg_case4_fY} then requires
\begin{equation*}
    f_2(X_2) = f_1(X_2) + (c_2 d_2 - c_1 d_1)X_2 = f_1(X_2) + \bigl(c_2(c_1 + d_1 - c_2) - c_1 d_1\bigr)X_2,
\end{equation*}
which is the expression in Equation~\eqref{eq:diffalg_case4_nonident_expr}. Since $c_2$ may be chosen freely, infinitely many valid triples $(f_2, c_2, d_2)$ exist, hence $f$ (as well as $c$ and $d$) is structurally non-identifiable. The identifiable quantities can be determined directly from Equation~\eqref{eq:diffalg_model2_Yobs_Yonly}. Here, any triplets $(f_1,c_1,d_1)$ and $(f_2,c_2,d_2)$ that preserve the quantities $c+d$ and $f(x) - cdx$ yield identical input-output equations, preventing us from identifying more detailed dynamics.
\end{proof}

We note that the non-identifiability relationship described in Equation~\eqref{eq:diffalg_case4_nonident_expr} only exhibits entanglement between the unknown function and a linear term, i.e. all higher order dynamics can still be identified. This is the same as in Figure~\ref{fig:single_var_nonident}, where the entanglement was also limited to the linear dynamics.

Case (ii) of Proposition~\ref{prop:diffalg_two_parameter_identifiability} is the only structurally non-identifiable case among the four mutual-activation-loop examples considered in this section, which illustrates that functional identifiability depends on both the model structure and the choice of observable. The non-identifiability here can, nevertheless, be resolved by incorporating additional information. For example, if only $X_2$ is observed but the initial conditions of both $X_1$ and $X_2$ are known for each observed trajectory, then all unknown functions and parameters become identifiable for generic initial conditions (proof in Appendix~\ref{appendix:proof_case4_known_initial_conditions}). 

% ------------------------------------------------------------------------------------------------------------------------------------ %
% ------------------------------ Results 3: Different forms of structural functional non-identifiability ----------------------------- %
% ------------------------------------------------------------------------------------------------------------------------------------ %

\section{Different forms of structural functional non-identifiability} \label{sec:nonidentifiability_types}

In classical structural parameter identifiability, non-identifiability arises when multiple parameter combinations generate identical observable dynamics. Proposition~\ref{prop:diffalg_two_parameter_identifiability}(ii) showed that this idea extends naturally to entanglement between unknown functions and parameters, where changes in one can be compensated by changes in the other. Here, we demonstrate that functional identifiability introduces further non-identifiability mechanisms that are absent from the classical parametric setting. We also revisit the distinction between local and global identifiability for unknown functions. A proof for each proposition is given in Appendix~\ref{appendix:proofs_section4}.

% --------- Function-to-function non-identifiability --------- %
\subsection{Function-to-function non-identifiability} \label{sec:nonidentifiability_types_function2function}

Consider the model
\begin{align}
    \dot{X}_1 &= f(X_2) + dX_1, \label{eq:f2f_nonident_model_dXdt} \\
    \dot{X}_2 &= X_1 - g(X_2), \label{eq:f2f_nonident_model_dYdt}
\end{align}
where $f:\mathbb{R}\to\mathbb{R}$ and $g:\mathbb{R}\to\mathbb{R}$ are unknown functions and $d\in\mathbb{R}$ is an unknown parameter.

\begin{proposition} \label{prop:f2f_nonident}
If only $X_2$ is observable, then the unknown functions $f$ and $g$ are both structurally non-identifiable. Specifically, given any admissible triplet $(f_1, g_1, d_1)$, any alternative choice of the form
\begin{align}
    \label{eq:f2f_nonident_expr}
    d_2&\in\mathbb{R},  \notag \\
    g_2(X_2) &= g_1(X_2) + (d_2 - d_1)X_2, \notag \\
    f_2(X_2) &= f_1(X_2) + d_1 g_1(X_2) - d_2 g_1(X_2) - d_2(d_2 - d_1)X_2,
\end{align}
yields dynamics identical to $(f_1, g_1, d_1)$ with respect to $X_2$.
\end{proposition}

Here we have a non-identifiability relationship involving multiple functions: a change in $g$ propagates to a compensating change in $f$ that preserves the dynamics. While not surprising, this shows how, while we previously only had non-identifiability relationships involving parameters, with the introduction of unknown functions, we can have relations with any combination of parameters and/or functions.

% --------- Pure functional non-identifiability --------- %

\subsection{Intrinsic functional non-identifiability} \label{sec:nonidentifiability_types_pure}

The previous examples all involved an unknown function co-occurring with other unknowns. We now show that a function can be non-identifiable even in the absence of any other unknown quantities. Consider the following model with a sole unknown, $f:\mathbb{R}\to\mathbb{R}$:
\begin{align}
    \dot{X}_1 &= f(X_2) - X_1,   \label{eq:fonly_nonident_model_dXdt} \\
    \dot{X}_2 &= (X_1 - 1)X_2.   \label{eq:fonly_nonident_model_dYdt}
\end{align}

\begin{proposition} \label{prop:pure_nonident}
If only $X_1$ is observed, then the unknown function $f$ is structurally non-identifiable due to a scaling non-identifiability. Specifically, given any admissible function $f_1$, any alternative function of the form
\begin{equation} \label{eq:pure_nonident_expr}
    f_2(X_2) = f_1(X_2/k), \qquad k\in\mathbb{R}\setminus\{0\},
\end{equation}
yields identical dynamics to those of $f_1$ with respect to $X_1$, provided $f$ is locally invertible.
\end{proposition}

This demonstrates another form of structural functional identifiability, further illustrating how functional identifiability introduces new nuance to the concept of identifiability.

% --------- Global and local structural identifiability --------- %

\subsection{Global and local structural identifiability} \label{sec:local_global}
In Definition \ref{def:functional_identifiability}, we noted that identifiability can be classified as either \emph{global} or \emph{local}. This distinction is well-established for classical structural parameter identifiability, an example of which can be found in the following self-activating gene regulatory network
\begin{equation*}
    \dot{X} = \frac{X^2}{X^2 + K^2} - dX.
\end{equation*}
Here, $d$ is globally identifiable, whereas $K$ is only locally identifiable. Specifically, the quantity $K^2$ is uniquely determined by the dynamics, but one cannot distinguish between $K$ and $-K$, both of which produce identical dynamics. Additional constraints, such as restricting to $K>0$, can restore global identifiability.

The same global/local distinction applies to functions. Consider the model
\begin{equation}
    \label{eq:local_global_model}
    \dot{X} = f(X) + \frac{1}{f(X)},
\end{equation}
where $f:\mathbb{R}\to\mathbb{R}$ is an unknown function.

\begin{proposition} \label{prop:local_global}
The unknown function $f$ is locally but not globally structurally identifiable. Specifically, given any admissible function $f_1$, the alternative function
\begin{equation} \label{eq:local_global_nonident}
    f_2(X) = \frac{1}{f_1(X)},
\end{equation}
yields identical dynamics, so $f$ is determined only up to the two-element equivalence class $\{f,\; 1/f
\}$.
\end{proposition}

Other local identifiability functional ambiguities can arise from sign symmetries, argument symmetries, and label-swapping between multiple functions; examples of which are given in Appendix~\ref{appendix:local_identifiability_examples}.

% ------------------------------------------------------------------------------------------------------------------------------------ %
% -------------------------------------------- Results 4: Applications to relevant models -------------------------------------------- %
% ------------------------------------------------------------------------------------------------------------------------------------ %

\section{Applications} \label{sec:applications}

We now apply the framework to two model classes with direct real-world relevance, a chemical reaction network model and the Lotka--Volterra model. In each case, we describe when the model exhibits structural identifiability and non-identifiability, and the form of any non-identifiabilities. Proofs are given in Appendix~\ref{appendix:proofs_section5}.

% --------- Chemical reaction networks --------- %

\subsection{Chemical reaction networks} \label{sec:applications_chemical reaction network}

Chemical reaction networks are a widely used modelling framework in biology and chemistry~\cite{hahl_comparison_2016,feinberg_foundations_2019}, with applications ranging from systems biology~\cite{goodwin_oscillatory_1965} and epidemiology~\cite{avram_advancing_2024} to pharmacology~\cite{zou_application_2020} and chemical kinetics~\cite{field_oscillations_1974,warnatz_combustion_1999}. Ordinary differential equation models are typically derived from reaction networks via the law of mass action. Importantly, the resulting systems often contain parameters or functional relationships that appear in multiple equations. Such repeated occurrences introduce additional constraints that can substantially affect identifiability. In this section, we use a chemical reaction network model both to illustrate structural functional identifiability analysis for this important class of systems and to demonstrate how identifiability is influenced when an unknown function appears across multiple equations.

We consider a simple two-species network with ODE model given by 
\begin{align}
    \dot{X}_1 &= f(X_2) - dX_1, \label{eq:chemical reaction network_dXdt} \\
    \dot{X}_2 &= X_1 - f(X_2), \label{eq:chemical reaction network_dYdt}
\end{align}
where $f:\mathbb{R}\to\mathbb{R}$ is an unknown function and $d\in\mathbb{R}$ is an unknown parameter. 

\begin{proposition} \label{prop:chemical reaction network}
When only $X_2$ is observable, $f$ and $d$ are structurally identifiable, given $d \neq 1$ and $f$ is non-affine. If $d = 1$, then only $f'$, but not $f$, is identifiable. When $f$ is affine, $f$ and $d$ are locally identifiable.
\end{proposition}

When $f$ is affine, identifiability is only local: writing $f(x) = ax + b$, the combinations $a + d$ and $(1-d)a$ are globally identifiable, but $a$ and $d$ individually admit two solutions (the unordered pair $\{a, d\}$ is recoverable, but their permutation is not). Full identifiability, therefore, requires a non-affine $f$.

\begin{proposition} \label{prop:chemical reaction network_Xobs}
When only $X_1$ is observable, $f$ is structurally non-identifiable but $d$ is structurally identifiable, provided $f$ is locally invertible. Specifically, given any admissible triplet $(f_1, d, X_{2,1}(0))$, any alternative choice of the form
\begin{align}
    \label{eq:chemical reaction network_Xobs_nonident_expr}
    k&\in\mathbb{R}, \notag \\
    f_2(X_2) &= f_1(X_2 + k), \notag \\
    X_{2,2}(0) &= X_{2,1}(0) - k,
\end{align}
yields dynamics identical to $(f_1, d, X_{2,1}(0))$ with respect to $X_1$. Here, $X_2(0)$ is the initial value of the unobserved variable $X_2$.
\end{proposition}

Again, this example demonstrates how identifiability depends on which variables are observable. It also demonstrates a case of intrinsic functional non-identifiability, i.e. that which does not involve entanglement with non-identifiable parameters. Finally, it illustrates how information regarding initial conditions is relevant: if we not only could observe $X_1$, but also the initial condition for  $X_2$, full identifiability would be achieved.

% --------- The Lotka--Volterra model --------- %

\subsection{The Lotka--Volterra model} \label{sec:applications_LV}

The Lotka--Volterra model is a classic example of predator--prey population dynamics~\cite{segel_j_2003}. The model has also seen extensive use to demonstrate hybrid modelling workflows, primarily for universal differential equations~\cite{dandekar_bayesian_2022,fronk_bayesian_2024,giampiccolo_robust_2024,grigorian_learning_2024,rackauckas_universal_2021}. Furthermore, it is an example where the unknown function depends on multiple variables, introducing additional complexity. We will consider two variants that differ in how the interaction term between the two species is represented. In each case, we will consider the case where either a single, or both, of the species are observable.

\paragraph{Model 1.} The two interaction terms are represented by distinct functions:
\begin{align}
    \dot{X}_1 &= \alpha X_1 + f(X_1,X_2), \label{eq:LV1_dXdt} \\
    \dot{X}_2 &= g(X_1,X_2) + \delta X_2. \label{eq:LV1_dYdt}
\end{align}
where $f:\mathbb{R}\to\mathbb{R}$ and $g:\mathbb{R}\to\mathbb{R}$ are unknown functions and $\alpha\in\mathbb{R}$ and $\delta\in\mathbb{R}$ are unknown parameters.

\begin{proposition} \label{prop:LV1_nonident}
For Model~1, $f$ and $g$ are structurally non-identifiable regardless of whether $X_1$, $X_2$, or both are observable.
\end{proposition}

As Model~1 is a fully augmented ODE of the form described in Section~\ref{sec:nonident_complete_UDEs}, this result follows directly from Proposition~\ref{prop:complete_UDE_nonident}.

\paragraph{Model 2.} In contrast to Model~1, the two interaction terms are assumed to be identical up to a linear rescaling:
\begin{align}
    \dot{X}_1 &= \alpha X_1 + f(X_1,X_2),        \label{eq:LV2_dXdt} \\
    \dot{X}_2 &= \gamma f(X_1,X_2) + \delta X_2, \label{eq:LV2_dYdt}
\end{align}
where $f:\mathbb{R}\to\mathbb{R}$ is an unknown function and $\alpha\in\mathbb{R}$, $\gamma\in\mathbb{R}$, and $\delta\in\mathbb{R}$ are unknown parameters.

\begin{proposition} \label{prop:LV2_observability}\label{prop:LV2_Xobs}\label{prop:LV2_remaining_observability}\label{prop:LV2_Yobs}\label{prop:LV2_both}
For Model~2, the following hold.
\begin{enumerate}
    \item When only $X_1$ is observable, $f$ is structurally non-identifiable. Specifically, given any admissible pair $(f_1, \alpha_1)$ and associated unobserved trajectory $X_{2,1}(t)$, any choice of the form
\begin{align}
    \label{eq:LV2_nonident_expr}
    \alpha_2&\in\mathbb{R}, \notag \\
    f_2(X_1, X_{2,2}) &= f_1(X_1, X_{2,1}) + (\alpha_1 - \alpha_2)X_1, \notag \\
    X_{2,2}(t)         &= X_{2,1}(t) + Z(t),
\end{align}
    where $Z$ solves $\dot{Z} = \gamma(\alpha_1 - \alpha_2)X_1 + \delta Z$, yields dynamics identical to $(f_1, \alpha_1, X_{2,1}(t))$.
    \item When only $X_2$ is observable, $f$ is structurally non-identifiable. Specifically, given any admissible pair $(f_1, \alpha, \gamma, \delta_1)$ and associated unobserved trajectory $X_{1,1}(t)$, any choice of the form
\begin{align} \label{eq:LV2_Yobs_nonident_expr}
    \delta_2&\in\mathbb{R}, \notag \\
    f_2(X_{1,2}, X_2) &= f_1(X_{1,1}, X_2) + \frac{(\delta_1 - \delta_2)X_2}{\gamma}, \notag \\
    X_{1,2}(t)      &= X_{1,1}(t) + Z(t),
\end{align}
    where $Z$ solves $\dot{Z} = \alpha Z + (\delta_1 - \delta_2)X_2/\gamma$, yields dynamics identical to $(f_1, \alpha, \gamma, \delta_1)$ with respect to $X_2$.
    \item When both $X_1$ and $X_2$ are observable, $f$, $\alpha$, $\gamma$, and $\delta$ are all structurally identifiable.
\end{enumerate}
\end{proposition}

The two single-observable cases in Proposition~\ref{prop:LV2_observability} differ from the preceding examples because their non-identifiability depends on the unobserved dynamics themselves. The ambiguity cannot be described solely in terms of transformations of the unknown functions and parameters; it also involves the unobserved state variable. More generally, unobserved states can be regarded as unknown quantities subject to identifiability analysis, alongside parameters and functions, although we do not pursue that perspective here. Case~(iii) is the only structurally functionally identifiable Lotka--Volterra case considered here, although alternative formulations can also achieve identifiability, for example by assuming certain parameters are known. Together, the examples highlight an important distinction between identifiability of model components and identifiability of the overall dynamics. Model~1, for instance, falls into the class of hybrid models discussed in Section~\ref{sec:nonident_neural_odes}, where the only identifiable expressions are system right-hand side expressions. As a consequence, this hybrid model formulation contributes no mechanistic knowledge, and is in practice equivalent to a pure Neural ODE.

% ------------------------------------------------------------------------------------------------------------------------------------ %
% ------------------------------------------------------------ Discussion ------------------------------------------------------------ %
% ------------------------------------------------------------------------------------------------------------------------------------ %

\section{Discussion} \label{sec:discussion}

In this work, we have extended the concept of structural identifiability from scalar parameters to unknown functions. Although functional and parametric identifiability are conceptually related, and can both be analysed using differential algebra, our results show that the functional setting introduces phenomena that have no direct analogue in the classical parameter identifiability framework. In addition to entanglement between unknown functions and parameters, we identified non-identifiability relationships involving function scaling, interactions between multiple unknown functions, and dependencies on unobserved system dynamics. These examples illustrate that moving from parameter inference to function inference is not simply a change in dimensionality, but introduces qualitatively new forms of ambiguity.

The transition from parameters to functions also affects the identifiability analysis itself. While the differential algebra approach remains applicable, the resulting derivations often require additional mathematical considerations that do not arise in the parametric setting. For example, the proof of Proposition~\ref{prop:diffalg_one_parameter_identifiability}(ii) relies on properties of the inverse of the unknown function, and several other examples require reasoning about functional transformations rather than parameter substitutions. Functional identifiability therefore extends not only the scope of structural identifiability analysis, but also the range of mathematical phenomena that such analyses must accommodate (like considering whether unknown functions are locally invertible or affine).

The framework demonstrated here is largely agnostic to the particular hybrid modelling methodology employed. It applies directly to approaches in which unknown functions are embedded within differential equations, including universal differential equations, physics-informed neural networks, and biology-informed neural networks. The framework can, in theory, also extend to equation discovery methods, such as SINDy or symbolic regression-based approaches. Here, candidate equations are viewed as functional relationships between state variables, and the validity of results requires that terms in non-identifiable functions' entanglement expressions also lie in the set of basis functions used for symbolic regression. More broadly, any modelling approach that seeks to infer unknown functional components within a fixed dynamical system falls within the scope of the theory developed here. By contrast, extending structural functional identifiability to settings in which the model structure itself changes, for example through the introduction of additional state variables and governing equations, remains an open challenge.

Although structural parameter identifiability is now a well-established field, its scope continues to expand. Recent work has extended identifiability analysis beyond ordinary differential equations to settings including partial differential equations and stochastic differential equations~\cite{renardy_structural_2022,browning_structural_2024,byrne_algebraic_2025,browning_exact_2025}. An important direction for future research is, therefore, to determine how functional identifiability can be formulated and analysed in these broader modelling frameworks. Further theoretical development is also needed to better accommodate more general classes of unknown functions, including functions of multiple variables, independent variables such as time, and unknown functions that depend on unknown parameters.

From a methodological perspective, applying differential algebra to unknown functions is often more involved than in the classical parametric setting. While the examples considered here remain tractable, it is not yet clear how well the approach scales to larger or more complex models. Furthermore, the widespread adoption of structural parameter identifiability has been facilitated by the development of automated software tools~\cite{structidjl}. Determining whether comparable automation is possible for structural functional identifiability, and developing the corresponding computational infrastructure, will be essential if functional identifiability analysis is to be applied routinely across scientific disciplines.

As scientific modelling becomes increasingly data-driven, a central question is what role mechanistic knowledge should play. Hybrid models offer a principled way to combine domain expertise with flexible function inference, but their flexibility also introduces the risk that multiple explanations may fit the same observations equally well. Structural functional identifiability provides a means of determining when mechanistic assumptions genuinely constrain the recovery of unknown dynamics and when they do not. In this sense, it offers a theoretical foundation for assessing whether hybrid models deliver new scientific understanding or simply alternative representations of the same observed behaviour.

% ------------------------------------------------------------------------------------------------------------------------------------ %
% --------------------------------------------------------- Acknowledgements --------------------------------------------------------- %
% ------------------------------------------------------------------------------------------------------------------------------------ %

\section*{Acknowledgements}

TEL and REB are supported by a grant from the Simons Foundation (MP-SIP-00001828). For the purpose of open access, the author has applied a CC BY public copyright licence to any author-accepted manuscript arising from this submission.

% ------------------------------------------------------------------------------------------------------------------------------------ %
% --------------------------------------------------------- Data availability -------------------------------------------------------- %
% ------------------------------------------------------------------------------------------------------------------------------------ %

% \section*{Code availability}

% All scripts for regenerating all figures (as well as all other numerical analyses) can be found here: \texttt{[GitHub link to be added]}.

% ------------------------------------------------------------------------------------------------------------------------------------ %
% ----------------------------------------------------------- Bibliography ----------------------------------------------------------- %
% ------------------------------------------------------------------------------------------------------------------------------------ %

\bibliographystyle{RS}
\bibliography{references}

% ------------------------------------------------------------------------------------------------------------------------------------ %
% ------------------------------------------------------------- Appendix ------------------------------------------------------------- %
% ------------------------------------------------------------------------------------------------------------------------------------ %
\appendix

% ------------------------------------------------------------------------------------------------------------------------------------ %
% ----------------------------------------------------- Appendix A: All Proofs ------------------------------------------------------- %
% ------------------------------------------------------------------------------------------------------------------------------------ %

\newpage

\section{Proofs} \label{appendix:proofs}

% --------- Proofs for Section~\ref{sec:general_nonident} --------- %

\subsection{Proofs for Section~\ref{sec:general_nonident}} \label{appendix:proofs_section2}

\subsubsection{Proof of Proposition~\ref{prop:1var_nonident}} \label{appendix:proof_1var}

Given an admissible choice $(f_1,\vec{p}_1)$, we wish to show that any proposed choice of the form
\begin{align*}    
    \vec{p}_2 &\in \mathbb{R}^{|\vec{p}_2|}, \\
    f_2(X) &= f_1(X) + g(X;\vec{p}_1) - g(X;\vec{p}_2),
\end{align*}
yields identical dynamics. Inserting these into Equation~\eqref{eq:nonident_1var_UDEs} yields
\begin{align*}
    \dot{X} &= f_2(X) + g(X;\vec{p}_2) \\
            &= f_1(X) + g(X;\vec{p}_1) - g(X;\vec{p}_2) + g(X;\vec{p}_2) \\
            &= f_1(X) + g(X;\vec{p}_1).
\end{align*}
This is identical to the dynamics generated by $(f_1,\vec{p}_1)$. Since $\vec{p}_2$ may be chosen freely, infinitely many valid $f_2$ exist, hence $f$ (as well as $\vec{p}$) is structurally non-identifiable.

\subsubsection{Proof of Proposition~\ref{prop:multivar_nonident}} \label{appendix:proof_multivar}

Given an admissible choice with parameter split $\vec{p}_{1,1}=\{\vec{a}_1,\vec{b}\}$, consider any proposed choice with $\vec{p}_{1,2}=\{\vec{a}_2,\vec{b}\}$ of the form
\begin{align*}
    \vec{a}_2    &\in \mathbb{R}^{|\vec{a}|}, \\
    f_2(\vec{u}) &= f_1(\vec{u}) + r(\vec{v};\vec{a}_1) - r(\vec{v};\vec{a}_2),
\end{align*}
yields identical dynamics. Inserting these into Equation~\eqref{eq:nonident_partial_system_UDEs} yields
\begin{align*}
    \dot{X}_1 &= f_2(\vec{u}) + g_1(\vec{X};\vec{p}_{1,2}) \\
              &= f_2(\vec{u}) + r(\vec{v};\vec{a}_2) + s(\vec{X};\vec{b}) \\
              &= f_1(\vec{u}) + r(\vec{v};\vec{a}_1) - r(\vec{v};\vec{a}_2) + r(\vec{v};\vec{a}_2) + s(\vec{X};\vec{b}) \\
              &= f_1(\vec{u}) + r(\vec{v};\vec{a}_1) + s(\vec{X};\vec{b}) \\
              &= f_1(\vec{u}) + g_1(\vec{X};\vec{p}_{1,1}).
\end{align*}
This is identical to the dynamics generated by $(f_1, \vec{p}_{1,1})$. The remaining equations are unchanged, as neither $f$ nor $\vec{a}$ appears in these. Since $\vec{a}_2$ may be chosen freely, infinitely many valid choices of $f_2$ exist, hence $f$ is structurally non-identifiable.

\subsubsection{Repeated-function counter-proof} \label{appendix:repeated_function_counterexample}

The proof of Proposition~\ref{prop:multivar_nonident} requires that $f$ appears only in the first equation. To see why this matters, consider a repeated-function extension
\begin{align}
    \dot{X}_1 &= f(\vec{u}) + g_1(\vec{X};\vec{p}_1), \label{eq:repeated_f_counter_X1} \\
    \dot{X}_2 &= f(\vec{u}) + g_2(\vec{X};\vec{p}_2), \label{eq:repeated_f_counter_X2} \\
    \dot{X}_i &= g_i(\vec{X};\vec{p}_i), \qquad i=3,\ldots,n, \notag
\end{align}
where $f$ appears in both the first and second equation. Here, identifiability may still hold for $f$.

\begin{proof}
The construction in Proposition~\ref{prop:multivar_nonident} changes $\vec{a}_1$ to an arbitrary $\vec{a}_2$ and compensates in the first equation by setting
\begin{equation*}
    f_2(\vec{u}) = f_1(\vec{u}) + r(\vec{v};\vec{a}_1) - r(\vec{v};\vec{a}_2).
\end{equation*}
This still preserves Equation~\eqref{eq:repeated_f_counter_X1}, since
\begin{align*}
    f_2(\vec{u}) + r(\vec{v};\vec{a}_2) + s(\vec{X};\vec{b})
    = f_1(\vec{u}) + r(\vec{v};\vec{a}_1) + s(\vec{X};\vec{b}) = f_1(\vec{u}) + g_1(\vec{X};\vec{p}_1).
\end{align*}
However, the same substitution changes Equation~\eqref{eq:repeated_f_counter_X2}:
\begin{align*}
    f_2(\vec{u}) + g_2(\vec{X};\vec{p}_2)
    = f_1(\vec{u}) + r(\vec{v};\vec{a}_1) - r(\vec{v};\vec{a}_2) + g_2(\vec{X};\vec{p}_2).
\end{align*}
This is identical to the original second equation only if $r(\vec{v};\vec{a}_1)=r(\vec{v};\vec{a}_2)$ on the admissible domain. Thus the arbitrary choice of $\vec{a}_2$ that generates the non-identifiable family in Proposition~\ref{prop:multivar_nonident} is no longer available. A trivial case where $f$ remains identifiable is where $g_2$ depends on no unknown parameters, i.e.~$|\vec{p}_2| = 0$. In this case, Equation~\eqref{eq:repeated_f_counter_X2} directly determines $f(\vec{u})$ from the observed dynamics (without even needing to consider Equation~\eqref{eq:repeated_f_counter_X1}). This does not prove that every repeated-function model is identifiable, but it shows that Proposition~\ref{prop:multivar_nonident} cannot be extended unchanged once $f$ appears in another equation.
\end{proof}

\subsubsection{Proof of Proposition~\ref{prop:complete_UDE_nonident}} \label{appendix:proof_complete_UDE}

Given an admissible choice $(f_{i,1}, \vec{p}_{i,1})$, choose any consistent alternative parameter sets $\vec{p}_{i,2}$, respecting any overlaps between the $\vec{p}_i$, and define
\begin{align*}
    f_{i,2}(\vec{X}) &= f_{i,1}(\vec{X}) + g_i(\vec{X};\vec{p}_{i,1}) - g_i(\vec{X};\vec{p}_{i,2}).
\end{align*}
Inserting these into Equation~\eqref{eq:complete_UDE} gives
\begin{align*}
    \dot{X}_i &= f_{i,2}(\vec{X}) + g_i(\vec{X};\vec{p}_{i,2}) \\
              &= f_{i,1}(\vec{X}) + g_i(\vec{X};\vec{p}_{i,1}) - g_i(\vec{X};\vec{p}_{i,2}) + g_i(\vec{X};\vec{p}_{i,2}) \\
              &= f_{i,1}(\vec{X}) + g_i(\vec{X};\vec{p}_{i,1}).
\end{align*}
This is identical to the dynamics generated by $(f_{i,1}, \vec{p}_{i,1})$. Since the alternative parameter sets $\vec{p}_{i,2}$ may be varied consistently, infinitely many valid choices of $f_{i,2}$ exist, hence all $f_i$ (as well as all $\vec{p}_i$) are structurally non-identifiable.

\subsubsection{Proof of Proposition~\ref{prop:nonident_neural_odes}} \label{appendix:proof_neural_odes}

\textit{If all states are observed, the system is identifiable.} Since each $X_i$ is fully observed, both it and its differentials are fully known. From this, and the relation $\dot{X}_i = f_i(\vec{X})$, we have that each $f_i$ is fully identifiable.

\medskip\noindent
\textit{If at least one state is unobserved, the system is non-identifiable.} We will show that non-identifiability holds for the case where only a single variable, $X_1$, is unobservable. From this, it directly follows that non-identifiability still holds when observability for additional states is lost.

Take an admissible set of functions $f_{j,1}$ and an associated trajectory of the unobserved state $X_{1,1}(t)$. We wish to show that any proposed choice of the form
\begin{align*}
    k    &\in \mathbb{R}, \\
    f_{i,2}(X_1, X_2, ...) &= f_{i,1}(X_1 - k, X_2, \ldots),\\
    X_{1,2}(t) &= X_{1,1}(t) + k,
\end{align*}
yields identical dynamics. Inserting these into Equation~\eqref{eq:nonident_neural_odes} gives, for each $i$,
\begin{align*}
    \dot{X}_i &= f_{i,2}(X_{1,2}, X_2, \ldots) \\
              &= f_{i,1}(X_{1,2} - k, X_2, \ldots) \\
              &= f_{i,1}(X_{1,1} + k - k, X_2, \ldots) \\
              &= f_{i,1}(X_{1,1}, X_2, \ldots).
\end{align*}
This is identical to the dynamics generated by $(f_{i,1}, X_{1,1}(t))$. Since $k$ may be chosen freely, infinitely many valid choices of $f_{i,2}$ exist, hence all $f_i$ (as well as the dynamics of the unknown variable $X_1$) are structurally non-identifiable.

% --------- Proofs for Section~\ref{sec:diffalg} --------- %

\subsection{Proofs for Section~\ref{sec:diffalg}} \label{appendix:proofs_section3}

For Proposition~\ref{prop:diffalg_one_parameter_identifiability}, part~(i) is proved in the main text and part~(ii) below. For Proposition~\ref{prop:diffalg_two_parameter_identifiability}, part~(i) is proved below and part~(ii) in the main text.

\subsubsection{Proof of Proposition~\ref{prop:diffalg_one_parameter_identifiability}(ii)} \label{appendix:proof_case2}

Isolate $X_2$ from Equation~\eqref{eq:diffalg_model1_dXdt}, writing
\begin{equation*}
    X_2 = f^{-1}(\dot{X}_1 + dX_1).
\end{equation*}
Let $u = \dot{X}_1 + dX_1$ and $g = f^{-1}$, so that $X_2 = g(u)$ and $\dot{X}_2 = g'(u)\dot{u}$. Substituting into Equation~\eqref{eq:diffalg_model1_dYdt} and rearranging gives
\begin{equation}
    \label{eq:diffalg_model1_Xobs_Xeq}
    X_1 = g'(u)\dot{u} + dg(u).
\end{equation}
Suppose $(f_1,d_1)$ and $(f_2,d_2)$ yield identical $X_1$ dynamics. Writing $g_1 = f_1^{-1}$, $g_2 = f_2^{-1}$, $u_1 = \dot{X}_1 + d_1 X_1$, and $u_2 = \dot{X}_1 + d_2 X_1$ and inserting both into Equation~\eqref{eq:diffalg_model1_Xobs_Xeq}, expanding $\dot{u}_1 = \ddot{X}_1 + d_1\dot{X}_1$, $\dot{u}_2 = \ddot{X}_1 + d_2\dot{X}_1$, and rearranging gives:
\begin{align}
    \label{eq:diffalg_model1_Xobs_equiveq}
    g_1'(u_1)\dot{u_1} + d_1g_1(u_1) &= X_1 = g_2'(u_2)\dot{u_2} + d_2g_2(u_2) \notag \Longleftrightarrow \\
    (\ddot{X}_1 + d_1\dot{X}_1)g_1'(u_1) + d_1g_1(u_1) &= (\ddot{X}_1 + d_2\dot{X}_1)g_2'(u_2) + d_2g_2(u_2) \notag \Longleftrightarrow \\
    \ddot{X}_1\,g_1'(u_1) + d_1\bigl(\dot{X}_1g_1'(u_1) + g_1(u_1)\bigr) &=  \ddot{X}_1\,g_2'(u_2) + d_2\bigl(\dot{X}_1g_2'(u_2) + g_2(u_2)\bigr).
\end{align}
Here, $\ddot{X}_1$ can be varied independently of $X_1$ and $\dot{X}_1$, so the coefficients of $\ddot{X}_1$ in Equation~\eqref{eq:diffalg_model1_Xobs_equiveq} must agree. Hence $g_1'(u_1) = g_2'(u_2)$. From this, and Lemma~\ref{lemma:g_prime_equality} (shown below), we have that either $d_1 = d_2$ or both $g_1'$ and $g_2'$ are constant.

\medskip\noindent
\textit{Case 1: $d_1 = d_2 \equiv d$.} Then $u_1 = u_2 \equiv u$. Cancelling the $\ddot{X}_1$ terms from Equation~\eqref{eq:diffalg_model1_Xobs_equiveq} gives 
\begin{equation*}
    d(\dot{X}_1g_1'(u) + g_1(u)) = d(\dot{X}_1g_2'(u) + g_2(u)).
\end{equation*}
We already have $g_1' = g_2'$, and we thus have $g_1(u) = g_2(u)$ and hence $f_1 = f_2$. The issue of invertibility of $f$ is discussed in Remark~\ref{remark:proof_invertibility}.

\medskip\noindent
\textit{Case 2: $g_1'$ and $g_2'$ constant.} By Lemma~\ref{lemma:constant_g_prime} below, $d_1 = d_2$ and $g_1 = g_2$, hence $f_1 = f_2$.

\medskip\noindent
In both cases, $d_1 = d_2$ and $f_1 = f_2$, so both $d$ and $f$ are structurally identifiable (subject to the local invertibility and $d \neq 0$ conditions considered below).

\begin{remark} \label{remark:proof_invertibility}

Multiple proofs of identifiability throughout this work (specifically those where the unknown function depends on an unobserved variable) assume that $f$ can be inverted. Here, the identifiability arguments hold as long as $f$ is \emph{locally invertible} (i.e. does not contain any flat regions).
\begin{proof}
Local invertibility implies that $f' = 0$ only at isolated points. Hence, the domain of $f$ can be written as a countable union of maximal open intervals $\bigcup_k I_k$ on each of which $f$ is strictly monotone. On each $I_k$, a smooth local inverse $g_k\colon f(I_k) \to I_k$ exists. The main proof can be applied to each $I_k$ independently because the input-output equation is a pointwise relation: it constrains the values of $f$ and $f'$ at a point $y$, without coupling to the values of $f$ at any other point. Since $f|_{I_k}$ is invertible, the main proof establishes identifiability of $f$ on each $I_k$ and of all parameters. The parameter conclusions are constants, so they agree across intervals. The function identity $f_1 = f_2$ therefore holds on $\bigcup_k I_k$, which is dense in $\mathbb{R}$. Since $f_1 - f_2$ is continuous and vanishes on a dense set, $f_1 = f_2$ everywhere.
\end{proof}
\end{remark}

\begin{lemma} \label{lemma:g_prime_equality}
Let $u_1 = \dot{X}_1 + d_1 X_1$, $u_2 = \dot{X}_1 + d_2 X_1$ with $d_1 \neq 0$. If $g_1'(u_1) = g_2'(u_2)$ holds for all admissible values of $X_1$ and $\dot{X}_1$, then either $d_1 = d_2$ (and hence $u_1 = u_2$), or both $g_1'$ and $g_2'$ are constant functions.
\end{lemma}

\begin{proof}
If $d_1 = d_2$, then $u_1 = u_2$ and the conclusion holds trivially. Suppose therefore that $d_1 \neq d_2$; we show both $g_1'$ and $g_2'$ must be constant. Consider the line of potential values
\begin{align*}
    \dot{X}_1 &= s, \\
    X_1 &= -s/d_1,
\end{align*}
for $s\in\mathbb{R}$. Then
\begin{align*}
    u_1 &= s + d_1\!\left(-\frac{s}{d_1}\right) = 0, \\
    u_2 &= s + d_2\!\left(-\frac{s}{d_1}\right) = s\!\left(1 - \frac{d_2}{d_1}\right).
\end{align*}
Since $d_1 \neq d_2$, the factor $(1-d_2/d_1)$ is nonzero, so as $s$ is varied over $\mathbb{R}$, $u_2$ takes every value in $\mathbb{R}$. The hypothesis $g_1'(u_1) = g_2'(u_2)$ therefore gives $g_2'(w) = g_1'(0)$ for all $w \in \mathbb{R}$, so $g_2'$ is the constant $g_1'(0)$. Since $g_2'$ is constant,  $g_1'(u_1) = g_2'(u_2)$, and $u_1 = \dot{X}_1 + d_1 X_1$ can take any value in $\mathbb{R}$, it follows that $g_1'$ is constant as well, proving the lemma.
\end{proof}

\begin{lemma} \label{lemma:constant_g_prime}
Under the hypotheses of Lemma~\ref{lemma:g_prime_equality}, and assuming both pairs satisfy Equation~\eqref{eq:diffalg_model1_Xobs_equiveq}, if both $g_1'$ and $g_2'$ are constant, then $d_1 = d_2$ and $g_1 = g_2$.
\end{lemma}

\begin{proof}
Write $g_1(x) = ax + b_1$ and $g_2(x) = ax + b_2$ (the slopes are equal since $g_1' = g_2'$). Substituting into Equation~\eqref{eq:diffalg_model1_Xobs_equiveq} gives
\begin{align*}
    d_1(a\dot{X}_1 + a(\dot{X}_1 + d_1 X_1) + b_1) &= d_2(a\dot{X}_1 + a(\dot{X}_1 + d_2 X_1) + b_2).
\end{align*}
Expanding and collecting terms in $\dot{X}_1$, $X_1$, and constants:
\begin{equation*}
    2ad_1\dot{X}_1 + ad_1^2 X_1 + d_1 b_1 = 2ad_2\dot{X}_1 + ad_2^2 X_1 + d_2 b_2.
\end{equation*}
Equating coefficients gives $d_1 = d_2$ (from the $\dot{X}_1$ term, provided $a \neq 0$) and hence $b_1 = b_2$, so $g_1 = g_2$.
\end{proof}

\subsubsection{Proof of Proposition~\ref{prop:diffalg_two_parameter_identifiability}(i)} \label{appendix:proof_case3}

Isolate $X_2$ from Equation~\eqref{eq:diffalg_model2_dXdt} by writing
\begin{equation*}
    X_2 = f^{-1}(\dot{X}_1 + cX_1).
\end{equation*}
Set $u = \dot{X}_1 + cX_1$ and $g = f^{-1}$, so that $X_2 = g(u)$ and $\dot{X}_2 = g'(u)\dot{u}$ with $\dot{u} = \ddot{X}_1 + c\dot{X}_1$. Substituting into Equation~\eqref{eq:diffalg_model2_dYdt} and rearranging gives
\begin{equation}
    \label{eq:proofs_Xobs_f_ident_main}
    X_1 = g'(u)\dot{u} + dg(u).
\end{equation}
Suppose $(f_1,c_1,d_1)$ and $(f_2,c_2,d_2)$ yield identical $X_1$ dynamics. Write $g_1 = f_1^{-1}$, $g_2 = f_2^{-1}$, $u_1 = \dot{X}_1 + c_1 X_1$, and $u_2 = \dot{X}_1 + c_2 X_1$. Inserting both into~\eqref{eq:proofs_Xobs_f_ident_main}, expanding $\dot{u}_1 = \ddot{X}_1 + c_1\dot{X}_1$ and $\dot{u}_2 = \ddot{X}_1 + c_2\dot{X}_1$, and rearranging we have
\begin{align}
    \label{eq:proofs_Xobs_f_ident_equiv}
    g_1'(u_1)\dot{u}_1 + d_1 g_1(u_1) &= X_1 = g_2'(u_2)\dot{u}_2 + d_2 g_2(u_2) \notag \Longleftrightarrow \\
    (\ddot{X}_1 + c_1\dot{X}_1)g_1'(u_1) + d_1 g_1(u_1) &= (\ddot{X}_1 + c_2\dot{X}_1)g_2'(u_2) + d_2 g_2(u_2) \notag \Longleftrightarrow \\
    \ddot{X}_1\,g_1'(u_1) + c_1 g_1'(u_1)\dot{X}_1 + d_1 g_1(u_1) &= \ddot{X}_1\,g_2'(u_2) + c_2 g_2'(u_2)\dot{X}_1 + d_2 g_2(u_2).
\end{align}
Here, $\ddot{X}_1$ can be varied independently of $X_1$ and $\dot{X}_1$, so the coefficients of $\ddot{X}_1$ in Equation~\eqref{eq:proofs_Xobs_f_ident_equiv} must agree. Hence $g_1'(u_1) = g_2'(u_2)$. By Lemma~\ref{lemma:g_prime_equality} (with $c_1$, $c_2$ in place of $d_1$, $d_2$), either $c_1 = c_2$ or both $g_1'$ and $g_2'$ are constant (which implies that $g = f^{-1}$ is affine).

\medskip\noindent
\textit{Case 1: $c_1 = c_2 \equiv c$.} Then $u_1 = u_2 \equiv u$. Cancelling the $\ddot{X}_1$ and $\dot{X}_1$ terms from Equation~\eqref{eq:proofs_Xobs_f_ident_equiv} gives $d_1 g_1(u) = d_2 g_2(u)$. Differentiating gives $d_1 g_1'(u) = d_2 g_2'(u)$. Since $g_1' = g_2'$, either $d_1 = d_2$ (giving $g_1 = g_2$ and full identifiability) or $g_1' = g_2' = 0$ (so $g$ is constant, which contradicts local invertibility of $f$).

\medskip\noindent
\textit{Case 2: $g_1'$ and $g_2'$ constant.} Write $g_1(x) = ax + b_1$ and $g_2(x) = ax + b_2$ (since $g_1' = g_2'$). Separating terms in $\ddot{X}_1$, $\dot{X}_1$, $X_1$, and constants in Equation~\eqref{eq:proofs_Xobs_f_ident_equiv} gives
\begin{align*}
    c_1 + d_1 &= c_2 + d_2, \\
    c_1 d_1   &= c_2 d_2, \\
    b_1 d_1   &= b_2 d_2.
\end{align*}
The first two relations are symmetric with respect to $c$ and $d$. Hence we can only identify the two values of the set $\{c,d\}$, but not assign values to the individual parameters. The third relation then determines $b$ only up to the same two-fold ambiguity: if $\{c, d\} = {s_1, s_2}$, the two admissible intercepts are $b$ and $b\cdot s_1/s_2$. This yields the local identifiability equivalence classes of potential solutions $\{(c, d, f^{-1}(x) = ax + b) , (d, c, f^{-1}(x) = ax + bd/c)\}$.

\medskip\noindent
In Case~1, global identifiability holds. In Case~2, $f^{-1}$ must be affine, and local identifiability holds. Together, these cases establish part~(i) of Proposition~\ref{prop:diffalg_two_parameter_identifiability}.

\subsubsection{Identifiability with known initial conditions for the two-parameter mutual activation model} \label{appendix:proof_case4_known_initial_conditions}

\begin{proof}
Consider model~\eqref{eq:diffalg_model2_dXdt}--\eqref{eq:diffalg_model2_dYdt} with only $X_2$ observed, but suppose that $X_1(0)$ and $X_2(0)$ are known for each observed trajectory. Assume generic initial conditions, so that at least one trajectory has $X_2(0)\neq 0$. Since $X_2$ is observed, $\dot{X}_2(0)$ is also known. Evaluating Equation~\eqref{eq:diffalg_model2_dYdt} at $t=0$ gives
\begin{equation*}
    \dot{X}_2(0)=X_1(0)-dX_2(0),
\end{equation*}
and hence
\begin{equation*}
    d=\frac{X_1(0)-\dot{X}_2(0)}{X_2(0)}.
\end{equation*}
Thus $d$ is structurally identifiable.

Now suppose two choices $(f_1,c_1,d)$ and $(f_2,c_2,d)$ yield identical $X_2$ dynamics. Substituting both into the input-output equation~\eqref{eq:diffalg_model2_Yobs_Yonly} gives
\begin{equation*}
    -(c_1+d)\dot{X}_2 + f_1(X_2)-c_1dX_2
    =
    -(c_2+d)\dot{X}_2 + f_2(X_2)-c_2dX_2.
\end{equation*}
Since $\dot{X}_2$ and $X_2$ may be varied independently across generic trajectories, the coefficient of $\dot{X}_2$ gives $c_1+d=c_2+d$, and therefore $c_1=c_2$. Substituting this back into the same equation gives $f_1(X_2)=f_2(X_2)$. Hence $d$, $c$, and $f$ are all structurally identifiable.
\end{proof}

% --------- Proofs for Section~\ref{sec:nonidentifiability_types} --------- %

\subsection{Proofs for Section~\ref{sec:nonidentifiability_types}} \label{appendix:proofs_section4}

\subsubsection{Proof of Proposition~\ref{prop:f2f_nonident}} \label{appendix:proof_f2f}

Isolate $X_1$ from Equation~\eqref{eq:f2f_nonident_model_dYdt} and differentiate to give
\begin{align*}
    X_1       &= \dot{X}_2 + g(X_2), \\
    \dot{X}_1 &= \ddot{X}_2 + g'(X_2)\dot{X}_2.
\end{align*}
Substituting into Equation~\eqref{eq:f2f_nonident_model_dXdt} and rearranging we have
\begin{align}
    \label{eq:f2f_nonident_dY2dt2}
    \ddot{X}_2 + g'(X_2)\dot{X}_2 &= f(X_2) + d\bigl(\dot{X}_2 + g(X_2)\bigr) \notag \Longleftrightarrow \\
    \ddot{X}_2 &= \dot{X}_2(d - g'(X_2)) + f(X_2) + dg(X_2).
\end{align}
Suppose $(f_1, g_1, d_1)$ and $(f_2, g_2, d_2)$ yield identical $X_2$ dynamics. Selecting $f_2$, $g_2$, $d_2$ according to the forms in Equation~\eqref{eq:f2f_nonident_expr} and noting that this yields $g_2'(X_2) = g_1'(X_2) + (d_2 - d_1)$, we can insert into the right-hand side of Equation~\eqref{eq:f2f_nonident_dY2dt2} to give
\begin{align*}
    &\dot{X}_2(d_2 - g_2'(X_2)) + f_2(X_2) + d_2 g_2(X_2) \\
    & \qquad = \dot{X}_2(d_2 - g_1'(X_2) - (d_2 - d_1)) \\
    &\qquad \qquad  + f_1(X_2) + d_1 g_1(X_2) - d_2 g_1(X_2) - d_2(d_2-d_1)X_2 + d_2 g_1(X_2) + d_2(d_2-d_1)X_2 \\
    &\qquad = \dot{X}_2(d_1 - g_1'(X_2)) + f_1(X_2) + d_1 g_1(X_2).
\end{align*}
This expression equals the right-hand side of Equation~\eqref{eq:f2f_nonident_dY2dt2} for the first choice. Hence, the second-choice values yield identical $X_2$ dynamics. Since $d_2$ may be chosen freely and each choice produces distinct $f_2 \neq f_1$ and $g_2 \neq g_1$, both $f$ and $g$ are structurally non-identifiable.

\subsubsection{Proof of Proposition~\ref{prop:pure_nonident}} \label{appendix:proof_pure}

We isolate $X_2$ from Equation~\eqref{eq:fonly_nonident_model_dXdt} by writing
\begin{equation*}
    X_2 = f^{-1}(\dot{X}_1 + X_1).
\end{equation*}
Set $u = \dot{X}_1 + X_1$, so $X_2 = f^{-1}(u)$ and $\dot{X}_2 = (f^{-1})'(u)\,\dot{u}$. Substituting into Equation~\eqref{eq:fonly_nonident_model_dYdt} and rearranging gives
\begin{align}
    \label{eq:fonly_nonident_Xeq}
    (f^{-1})'(u)\,\dot{u} &= (X_1 - 1)\,f^{-1}(u) \notag \Longleftrightarrow \\
    X_1 &= \frac{(f^{-1}(u))'\dot{u}}{f^{-1}(u)} + 1.
\end{align}
Take $f_2$ as in Equation~\eqref{eq:pure_nonident_expr} for any $k \neq 0$. Then $f_2^{-1}(u) = k f_1^{-1}(u)$ and $(f_2^{-1})'(u) = k(f_1^{-1})'(u)$. Substituting gives
\begin{equation*}
    \frac{(f_2^{-1}(u))'\dot{u}}{f_2^{-1}(u)} + 1 = \frac{k(f_1^{-1}(u))'\dot{u}}{k f_1^{-1}(u)} + 1 = \frac{(f_1^{-1}(u))'\dot{u}}{f_1^{-1}(u)} + 1.
\end{equation*}
The first and second functions yield identical $X_1$ dynamics. Since $k$ may be chosen freely, $f$ is structurally non-identifiable.

\subsubsection{Proof of Proposition~\ref{prop:local_global}} \label{appendix:proof_local_global}

Substituting $f_2(X) = 1/f_1(X)$ into the right-hand side of Equation~\eqref{eq:local_global_model} gives
\begin{equation*}
    f_2(X) + \frac{1}{f_2(X)} = \frac{1}{f_1(X)} + \frac{1}{1/f_1(X)} = \frac{1}{f_1(X)} + f_1(X) = f_1(X) + \frac{1}{f_1(X)}.
\end{equation*}
The right-hand side is identical for $f_1$ and $f_2$, so both pairs yield identical $X$ dynamics. Since $f_2 = 1/f_1 \neq f_1$ whenever $f_1(X)^2 \neq 1$ for some $X$, the function $f$ is not globally identifiable. Finally, $(1/f_2)(X) = f_1(X)$, so the equivalence class contains exactly two elements, confirming local identifiability. Furthermore, even if $f$ is assumed to be continuous, whichever of the relations $f_1=f_2$ and $f_1=1/f_2$ holds may switch at points where $f_1(x^*)=f_2(x^*)$, permitting the generation of a potentially large equivalence class of functions.

% --------- Proofs for Section~\ref{sec:applications} --------- %

\subsection{Proofs for Section~\ref{sec:applications}} \label{appendix:proofs_section5}

\subsubsection{Proof of Proposition~\ref{prop:chemical reaction network}} \label{appendix:proof_chemical reaction network}

Isolate $X_1$ from Equation~\eqref{eq:chemical reaction network_dYdt} and differentiate to give
\begin{align*}
    X_1       &= \dot{X}_2 + f(X_2), \\
    \dot{X}_1 &= \ddot{X}_2 + f'(X_2)\dot{X}_2.
\end{align*}
Substituting into Equation~\eqref{eq:chemical reaction network_dXdt} we have
\begin{equation}
    \label{eq:dchemical reaction network_Yobs_dY2dt2}
    \ddot{X}_2 = (1-d)f(X_2) - (d + f'(X_2))\dot{X}_2.
\end{equation}
Suppose $(f_1,d_1)$ and $(f_2,d_2)$ yield identical $X_2$ dynamics. Inserting both into Equation~\eqref{eq:dchemical reaction network_Yobs_dY2dt2} and equating gives
\begin{equation}
    \label{eq:dchemical reaction network_Yobs_equiveq}
    (1-d_1)f_1(X_2) - (d_1 + f_1'(X_2))\dot{X}_2 = \ddot{X}_2 = (1-d_2)f_2(X_2) - (d_2 + f_2'(X_2))\dot{X}_2.
\end{equation}
Here, $\dot{X}_2$ and $X_2$ can be varied independently, allowing us to separate Equation~\eqref{eq:dchemical reaction network_Yobs_equiveq} into two independent equalities:
\begin{align}
    (1-d_1)f_1(X_2) &= (1-d_2)f_2(X_2), \label{eq:dchemical reaction network_Yobs_dalg_1} \\
    d_1 + f_1'(X_2) &= d_2 + f_2'(X_2). \label{eq:dchemical reaction network_Yobs_dalg_2}
\end{align}

\medskip\noindent
\textit{Case 1: $f$ is non-affine.} Suppose for contradiction that $d_1 \neq d_2$. From Equation~\eqref{eq:dchemical reaction network_Yobs_dalg_1}:
\begin{equation*}
    f_1(X_2) = \frac{1-d_2}{1-d_1} f_2(X_2).
\end{equation*}
Differentiating and substituting into Equation~\eqref{eq:dchemical reaction network_Yobs_dalg_2} gives $f_1'(X_2) = d_2 - 1$, a constant, so $f_1$ is affine. By the same argument applied to $f_2$, $f_2$ is also affine. This contradicts our assumption that $f$ is non-affine. Therefore $d_1 = d_2 \equiv d$. We then have:
\begin{itemize}
    \item if $d \neq 1$, Equation~\eqref{eq:dchemical reaction network_Yobs_dalg_1} gives $f_1 = f_2$, and so both $d$ and $f$ are globally structurally identifiable;
    \item if $d = 1$, Equation~\eqref{eq:dchemical reaction network_Yobs_dalg_1} is trivially satisfied for any $f_1$, $f_2$ and  Equation~\eqref{eq:dchemical reaction network_Yobs_dalg_2} gives $f_1' = f_2'$. Thus $d$ is identifiable, but only $f'$ (not $f$ itself) is identifiable.
\end{itemize}

\medskip\noindent
\textit{Case 2: $f$ is affine.} Write $f_i(X_2) = a_i X_2 + b_i$. Substituting into Equations~\eqref{eq:dchemical reaction network_Yobs_dalg_1}--\eqref{eq:dchemical reaction network_Yobs_dalg_2} and equating coefficients gives
\begin{align*}
    (1-d_1)a_1 &= (1-d_2)a_2, \\
    (1-d_1)b_1 &= (1-d_2)b_2, \\
    d_1 + a_1  &= d_2 + a_2.
\end{align*}
The first and third relations show that $a + d$ and $(1-d)a$ are globally identifiable. Substituting $a = (a+d) - d$ into $(1-d)a$ shows that $d$ satisfies the quadratic $d^2 - (1 + (a+d))d + ((a+d) - (1-d)a) = 0$, with coefficients determined by the globally identifiable quantities. This yields at most two solutions for $d$, and hence for $a$ and $b$. Thus $a$, $d$, and $b$ are each locally identifiable. 

In summary, full structural identifiability holds only in Case~1 with $d \neq 1$.

\subsubsection{Proof of Proposition~\ref{prop:chemical reaction network_Xobs}} \label{appendix:proof_chemical reaction network_Xobs}

Given an admissible choice $(f_1, d)$ with associated unobserved trajectory $X_{2,1}(t)$, we wish to show that the family in Equation~\eqref{eq:chemical reaction network_Xobs_nonident_expr} yields identical $X_1$ dynamics. Setting $X_{2,2}(t) = X_{2,1}(t) - k$, so that $X_{2,2}(0) = X_{2,1}(0) - k$ and $\dot{X}_{2,2} = \dot{X}_{2,1}$, and substituting $f_2(X_2) = f_1(X_2+k)$ into Equations~\eqref{eq:chemical reaction network_dXdt}--\eqref{eq:chemical reaction network_dYdt} gives
\begin{align*}
    \dot{X}_1 &= f_2(X_{2,2}) - dX_1 = f_1(X_{2,2} + k) - dX_1 = f_1(X_{2,1}) - dX_1, \\
    \dot{X}_{2,2} &= X_1 - f_2(X_{2,2}) = X_1 - f_1(X_{2,1}) = \dot{X}_{2,1}.
\end{align*}
Thus the second choice reproduces the first-choice $X_1$ dynamics exactly. Since $k$ may be chosen freely, $f$ is structurally non-identifiable.

We now show that $d$ is identifiable (provided $f$ is locally invertible). Set $g = f^{-1}$ and $u = \dot{X}_1 + dX_1$. From $\dot{X}_1 = f(X_2) - dX_1$ we have $X_2 = g(u)$, and differentiating gives $\dot{X}_2 = g'(u)\dot{u}$ with $\dot{u} = \ddot{X}_1 + d\dot{X}_1$. Substituting into Equation~\eqref{eq:chemical reaction network_dYdt} and rearranging yields
\begin{align} 
    g'(u)\dot{u} &= X_1 - f(g(u)) \Longleftrightarrow \notag \\
    g'(u)\dot{u} &= X_1 - u \Longleftrightarrow \notag \\
    g'(u)(\ddot{X}_1 + d\dot{X}_1) &= X_1 - (\dot{X}_1 + dX_1) \Longleftrightarrow \notag \\
    \dot{X}_1 &= (1-d)X_1 - g'(u)(\ddot{X}_1 + d\dot{X}_1). \label{eq:chemical reaction network_Xobs_Xdot}
\end{align}
Suppose $(g_1, d_1)$ and $(g_2, d_2)$ both yield identical $X_1$ dynamics, with $u_i = \dot{X}_1 + d_iX_1$. Inserting each into Equation~\eqref{eq:chemical reaction network_Xobs_Xdot}, equating (since both expressions equal the same $\dot{X}_1$), and rearranging, gives
\begin{align}
    (1-d_1)X_1 - g_1'(u_1)(\ddot{X}_1 + d_1\dot{X}_1) &= \dot{X}_1 = (1-d_2)X_1 - g_2'(u_2)(\ddot{X}_1 + d_2\dot{X}_1) \notag \Longleftrightarrow \\
    (1-d_1)X_1 -  d_1g_1'(u_1)\dot{X}_1 - g_1'(u_1)\ddot{X}_1 &= (1-d_2)X_1 -  d_2g_2'(u_2)\dot{X}_1 - g_2'(u_2)\ddot{X}_1. \label{eq:chemical reaction network_Xobs_equiveq}
\end{align}
Here, the observed variable and its derivatives can be varied independently, so the $X_1$ terms in Equation~\eqref{eq:chemical reaction network_Xobs_equiveq} must agree:
\begin{equation*}
    (1-d_1)X_1 = (1-d_2)X_1, 
\end{equation*}
which trivially gives $d_1 = d_2$, i.e. $d$ is structurally identifiable.

\subsubsection{Proofs for the Lotka--Volterra model} \label{appendix:proof_LV}

\begin{proof}[Proof of Proposition~\ref{prop:LV2_observability}(i)]
Given an admissible choice $(f_1, \alpha_1)$ with associated unobserved trajectory $X_{2,1}(t)$, we wish to show that any choice of the following form
\begin{align*}
    \alpha_2 &\in \mathbb{R}, \\
    f_2(X_1, X_{2,2}) &= f_1(X_1, X_{2,1}) + (\alpha_1 - \alpha_2)X_1,\\
    X_{2,2}(t)         &= X_{2,1}(t) + Z(t),
\end{align*}
where $Z$ solves $\dot{Z} = \gamma(\alpha_1 - \alpha_2)X_1 + \delta Z$, yields dynamics identical to $(f_1, \alpha_1, X_{2,1}(t))$. Inserting $(f_2, \alpha_2, X_{2,2}(t))$ into Equation~\eqref{eq:LV2_dXdt} gives
\begin{align}
    \dot{X}_1 &= \alpha_2 X_1 + f_2(X_1, X_{2,2}) \\
            &= \alpha_2 X_1 + f_1(X_1, X_{2,1}) + (\alpha_1 - \alpha_2)X_1 \\
            &= \alpha_1 X_1 + f_1(X_1,X_{2,1}).
\end{align}
The $X_1$ dynamics of $(f_1, \alpha_1, X_{2,1}(t))$ are therefore reproduced. It remains to verify that $X_{2,2} = X_{2,1} + Z$ satisfies the $\dot{X}_2$ equation~\eqref{eq:LV2_dYdt} for the second model. Since $X_{2,1}$ satisfies the first system and $Z$ satisfies $\dot{Z} = \gamma(\alpha_1 - \alpha_2)X_1 + \delta Z$ by construction we have
\begin{align*}
    \dot{X}_{2,2} &= \dot{X}_{2,1} + \dot{Z} \\
              &= \bigl(\gamma f_1(X_1,X_{2,1}) + \delta X_{2,1}\bigr) + \bigl(\gamma(\alpha_1 - \alpha_2)X_1 + \delta Z\bigr) \\
              &= \gamma\bigl(f_1(X_1,X_{2,1}) + (\alpha_1 - \alpha_2)X_1\bigr) + \delta(X_{2,1} + Z) \\
              &= \gamma f_2(X_1,X_{2,2}) + \delta X_{2,2}.
\end{align*}
Thus $(f_2, \alpha_2, X_{2,2})$ is a valid solution to the full system that reproduces the same $X_1$ dynamics. Since $\alpha_2$ may be chosen freely, infinitely many valid pairs $(f_2, X_{2,2})$ exist, so $f$ is structurally non-identifiable.
\end{proof}

\begin{proof}[Proof of Proposition~\ref{prop:LV2_observability}(ii)]
Given an admissible choice $(f_1, \alpha, \gamma, \delta_1)$ with associated trajectories $X_{1,1}(t)$ and $X_2(t)$, we wish to show that any choice of the form
\begin{align*}
    \delta_2    &\in \mathbb{R}, \\
    f_2(X_{1,2}, X_2) &= f_1(X_{1,1}, X_2) + \frac{(\delta_1 - \delta_2)X_2}{\gamma}, \\
    X_{1,2}(t)      &= X_{1,1}(t) + Z(t),
\end{align*}
where $Z$ solves $\dot{Z} = \alpha Z + (\delta_1 - \delta_2)X_2/\gamma$, yields dynamics identical to $(f_1, \alpha, \gamma, \delta_1)$ with respect to $X_2$. Inserting $(f_2, \delta_2, X_{1,2})$ into Equation~\eqref{eq:LV2_dYdt} gives
\begin{align*}
    \dot{X}_2 &= \gamma f_2(X_{1,2}, X_2) + \delta_2 X_2 \\
            &= \gamma\!\left(f_1(X_{1,1}, X_2) + \frac{(\delta_1 - \delta_2)X_2}{\gamma}\right) + \delta_2 X_2 \\
            &= \gamma f_1(X_{1,1}, X_2) + (\delta_1 - \delta_2)X_2 + \delta_2 X_2 \\
            &= \gamma f_1(X_{1,1}, X_2) + \delta_1 X_2.
\end{align*}
The $X_2$ dynamics of $(f_1, \alpha, \gamma, \delta_1)$ are therefore reproduced. It remains to verify that $X_{1,2} = X_{1,1} + Z$ satisfies Equation~\eqref{eq:LV2_dXdt} for the second model. Since $X_{1,1}$ satisfies the first system and $Z$ satisfies $\dot{Z} = \alpha Z + (\delta_1 - \delta_2)X_2/\gamma$ by construction we have
\begin{align*}
    \dot{X}_{1,2} &= \dot{X}_{1,1} + \dot{Z} \\
              &= \bigl(\alpha X_{1,1} + f_1(X_{1,1}, X_2)\bigr) + \left(\alpha Z + \frac{(\delta_1 - \delta_2)X_2}{\gamma}\right) \\
              &= \alpha(X_{1,1} + Z) + f_1(X_{1,1}, X_2) + \frac{(\delta_1 - \delta_2)X_2}{\gamma} \\
              &= \alpha X_{1,2} + f_2(X_{1,2}, X_2).
\end{align*}
Thus $(f_2, \alpha, \gamma, \delta_2, X_{1,2})$ is a valid solution to the full system that reproduces the same $X_2$ dynamics. Since $\delta_2$ may be chosen freely, infinitely many valid pairs $(f_2, X_{1,2})$ exist, so $f$ is structurally non-identifiable.
\end{proof}

\begin{proof}[Proof of Proposition~\ref{prop:LV2_observability}(iii)]
Suppose $(f_1, \alpha_1, \gamma_1, \delta_1)$ and $(f_2, \alpha_2, \gamma_2, \delta_2)$ yield identical $X_1$ and $X_2$ dynamics. Rewriting~\eqref{eq:LV2_dYdt} as $\dot{X}_2/\gamma = f(X_1,X_2) + (\delta/\gamma)X_2$ and subtracting from Equation~\eqref{eq:LV2_dXdt} gives
\begin{equation*}
    \dot{X}_1 - \dot{X}_2/\gamma = \alpha X_1 - (\delta/\gamma)X_2.
\end{equation*}
Since $X_1$, $X_2$, $\dot{X}_1$, and $\dot{X}_2$ are all observed and may be varied independently, $\alpha$, $\delta/\gamma$, and $1/\gamma$ are directly identifiable as the coefficients of $X_1$, $X_2$, and $\dot{X}_2$ in this expression, giving $\alpha_1 = \alpha_2$, $\gamma_1 = \gamma_2$, and $\delta_1 = \delta_2$ (provided $\gamma \neq 0$). With $\alpha$ now known, $f(X_1,X_2) = \dot{X}_1 - \alpha X_1$ is directly computable from observations, hence $f_1 = f_2$. Since there is only a single choice $(f, \alpha, \gamma, \delta)$ yielding any given dynamics, all functions and parameters are structurally identifiable.
\end{proof}

% ------------------------------------------------------------------------------------------------------------------------------------ %
% ---------------------------------------------------- Appendix B: Example models ----------------------------------------------------- %
% ------------------------------------------------------------------------------------------------------------------------------------ %

\section{Example models for Section~\ref{sec:general_nonident}} \label{appendix:example_models}

Below are showcase example models satisfying the conditions of each proposition in Section~\ref{sec:general_nonident}.

% --------- Prop 3.1 examples --------- %

\subsection{Proposition~\ref{prop:1var_nonident}: single-variable ODEs} \label{appendix:examples_1var}

Such models are of the form $\dot{X} = f(X) + g(X;\vec{p})$ with $f$ unknown and at least one unknown parameter in $g$.

\begin{enumerate}
    
    \item \textbf{Falling object with unknown drag.} Velocity $v$ of an object subject to gravity and an unknown drag force.
    \begin{equation*}
        \dot{v} = g_0 - f(v).
    \end{equation*}
    Here $-f(v)$ is the unknown drag law, describing how air resistance depends on speed and $g(v;g_0) = g_0$ is the known gravitational term, with unknown acceleration $g_0$.

    \item \textbf{Population growth with unknown density-dependent mortality.} A population of size $N$ growing linearly, subject to an unknown density-dependent mortality capturing effects such as crowding or resource competition.
    \begin{equation*}
        \dot{N} = rN - f(N).
    \end{equation*}
    Here $-f(N)$ is the unknown density-dependent mortality law and $g(N;r) = rN$ is the known linear growth term, with unknown per-capita rate $r$.

    \item \textbf{Tumour growth under treatment.} Tumour volume $V$ with an unknown growth law and linear treatment-induced clearance.
    \begin{equation*}
        \dot{V} = f(V) - kV.
    \end{equation*}
    Here $f(V)$ is the unknown growth or proliferation law and $g(V;k) = -kV$ is the known linear treatment-clearance term, with unknown rate $k$.

    \item \textbf{Newton's law of cooling with unknown radiation.} Temperature $T$ of an object cooling via convection to an ambient environment and via an unknown radiative process.
    \begin{equation*}
        \dot{T} = f(T) - h(T - T_{\mathrm{env}}).
    \end{equation*}
    Here $f(T)$ is the unknown radiative heat-loss law. $g(T;h,T_{\mathrm{env}}) = -h(T - T_{\mathrm{env}})$ is the known convective term, with unknown heat-transfer coefficient $h$ and unknown ambient temperature $T_{\mathrm{env}}$.
    
\end{enumerate}

% --------- Prop 3.2 examples --------- %

\subsection{Proposition~\ref{prop:multivar_nonident}: multi-variable ODEs} \label{appendix:examples_multivar}

 We consider the class of models where the equation containing $f$ decomposes as $\dot{X}_1 = f(\vec{u}) + r(\vec{v};\vec{a}) + s(\vec{X};\vec{b})$ with $\vec{v} \subseteq \vec{u}$ and $\vec{a}$ absent from all other equations. The key structural requirement is that the known part of the equation containing $f$ includes a parameter-dependent term $r(\vec{v};\vec{a})$ acting on variables $\vec{v}$ that are also arguments of $f$. Changes in $\vec{a}$ can always be absorbed into $f$ while leaving the dynamics unchanged, making both non-identifiable.

\begin{enumerate}

    \item \textbf{Self-activating gene-expression network.} Two biochemical species $X_1$ and $X_2$ interconverting at known rates, where $X_1$ additionally exhibits unknown self-activated production and both species degrade linearly.
    \begin{align*}
        \dot{X}_1 &= f(X_1) - k_1 X_1 + k_2 X_2 - d_1 X_1, \\
        \dot{X}_2 &= k_1 X_1 - k_2 X_2 - d_2 X_2.
    \end{align*}
    Here $f(\vec{u}) = f(X_1)$ is the unknown self-activation rate, with $\vec{u} = \{X_1\}$, $r(\vec{v};\vec{a}) = -d_1 X_1$ has $\vec{a} = \{d_1\}$ (the unknown degradation rate of $X_1$) acting on $\vec{v} = \{X_1\} \subseteq \vec{u}$, absent from $\dot{X}_2$. The function $s(\vec{X}) = -k_1 X_1 + k_2 X_2$ contains the known interconversion rates $k_1$ and $k_2$---since these are known (not inferred) they are not entangled with $f$, and $d_2$ is the unknown degradation rate of $X_2$.

    \item \textbf{Epidemic model with unknown nonlinear recovery.} Susceptible $S$ and infectious $I$ populations, where infectious individuals recover through both a linear recovery rate and a nonlinear rate capturing e.g.\ hospital-capacity constraints.
    \begin{align*}
        \dot{S} &= -\beta SI, \\
        \dot{I} &= \beta SI - f(I) - \gamma I.
    \end{align*}
    Here $f(\vec{u}) = f(I)$ is the unknown nonlinear recovery rate, with $\vec{u} = \{I\}$. $r(\vec{v};\vec{a}) = -\gamma I$ has $\vec{a} = \{\gamma\}$ (the unknown linear recovery rate) acting on $\vec{v} = \{I\} \subseteq \vec{u}$, absent from $\dot{S}$. The infection term, $s(\vec{X};\vec{b}) = \beta SI$ has $\vec{b} = \{\beta\}$ (the unknown transmission rate parameter).

    \item \textbf{Grazer and vegetation dynamics.} A grazer population $G$ consuming vegetation $V$, growing at an unknown rate depending on both $G$ and $V$. $G$ is also subject to constant immigration and linear mortality, while $V$ follows logistic growth.
    \begin{align*}
        \dot{G} &= f(G,V) + a - mG, \\
        \dot{V} &= rV\!\left(1 - \frac{V}{K}\right).
    \end{align*}
    Here $f(\vec{u}) = f(G,V)$ is the unknown growth law, with $\vec{u} = \{G,V\}$. $r(\vec{v};\vec{a}) = a - mG$ has $\vec{a} = \{a,m\}$ (the unknown immigration rate and mortality rate) acting on $\vec{v} = \{G\} \subseteq \vec{u}$, absent from $\dot{V}$; $s(\vec{X};\vec{b}) \equiv 0$ is trivial. The parameters $r$ and $K$ are the vegetation's unknown intrinsic growth rate and carrying capacity.
    
\end{enumerate}

% --------- Prop 3.3 examples --------- %

\subsection{Proposition~\ref{prop:complete_UDE_nonident}: Fully augmented ODEs} \label{appendix:examples_complete}

We consider the class of models where every equation has the form $\dot{X}_i = f_i(\vec{X}) + g_i(\vec{X};\vec{p}_i)$ with at least one unknown parameter per equation.

\begin{enumerate}
    \item \textbf{Coupled thermal compartments with unknown exchange laws.} Temperatures $T_1$ and $T_2$ in two coupled compartments, with unknown heat-exchange corrections and explicit linear heat loss to an ambient environment.
    \begin{align*}
        \dot{T}_1 &= f(T_1,T_2) - h_1(T_1 - T_{\mathrm{env}}), \\
        \dot{T}_2 &= g(T_1,T_2) - h_2(T_2 - T_{\mathrm{env}}).
    \end{align*}
    Here $f(T_1,T_2)$ and $g(T_1,T_2)$ are the unknown exchange laws. The known terms contain unknown heat-loss coefficients $h_1$ and $h_2$, which can be absorbed into the corresponding unknown functions.

    \item \textbf{Two competing firms with unknown market effects.} Revenue levels $x_1$ and $x_2$ of two firms operating in partially overlapping but distinct market segments, so that each firm's growth is limited both by its own market saturation and by competitive pressure from the other. Unknown functions $f$ and $g$ capture additional market dynamics not represented by the explicit competition structure---for instance, brand-loyalty effects or differential responses to external demand shocks that depend on both firms' current positions.
    \begin{align*}
        \dot{x}_1 &= f(x_1,x_2) + r_1 x_1\!\left(1 - \frac{x_1 + \alpha_{12}\,x_2}{K_1}\right), \\
        \dot{x}_2 &= g(x_1,x_2) + r_2 x_2\!\left(1 - \frac{x_2 + \alpha_{21}\,x_1}{K_2}\right).
    \end{align*}
    Here $f(x_1,x_2)$ and $g(x_1,x_2)$ are the unknown market-effect laws: $g_1(\vec{X};r_1,K_1,\alpha_{12}) = r_1 x_1(1-(x_1+\alpha_{12}x_2)/K_1)$ contains unknown intrinsic growth rate $r_1$, market capacity $K_1$, and cross-competition coefficient $\alpha_{12}$ (with $\alpha_{12} \neq \alpha_{21}$ reflecting the asymmetry of the two segments), while $g_2(\vec{X};r_2,K_2,\alpha_{21}) = r_2 x_2(1-(x_2+\alpha_{21}x_1)/K_2)$ contains the corresponding parameters $r_2$, $K_2$, and $\alpha_{21}$.

    \item \textbf{Soil moisture and vegetation water dynamics.} Soil moisture $M$ and vegetation water status $V$ (e.g.\ plant water content), coupled through unknown soil-water and plant-water dynamics laws, with a constant water input and linear loss terms.
    \begin{align*}
        \dot{M} &= f(M,V) + P - \ell M, \\
        \dot{V} &= g(M,V) - mV.
    \end{align*}
    Here $f(M,V)$ is the unknown soil--water dynamics law, capturing the effect of soil and vegetation water content on soil--water retention, and $g(M,V)$ is the unknown plant--water dynamics law, capturing the effect of soil moisture on vegetation water uptake. The function $g_1(\vec{X};P,\ell) = P - \ell M$ contains unknown constant water-input rate $P$ (e.g.\ average rainfall) and unknown background moisture-loss rate $\ell$, while $g_2(\vec{X};m) = -mV$ contains unknown vegetation-water loss rate $m$.
\end{enumerate}

% ------------------------------------------------------------------------------------------------------------------------------------ %
% ------------------------------------------- Appendix C: local functional identifiability examples ------------------------------------ %
% ------------------------------------------------------------------------------------------------------------------------------------ %

\section{Additional examples of local functional identifiability} \label{appendix:local_identifiability_examples}

In Section~\ref{sec:local_global} we introduced the notion of local functional identifiability, and gave a simple example in Proposition~\ref{prop:local_global}. Here, we give additional examples of different forms of local structural identifiability that can appear in the context of functions. The first two examples are translation from cases of local structural parametric identifiability, while others two only appears in the context of functions. These examples are not intended to be exhaustive, but rather to illustrate how a wide variety of local functional identifiability scenarios are possible.

\begin{enumerate}
    \item \textbf{Sign ambiguity.} Consider
    \begin{equation*}
        \dot{X} = f(X)^2,
    \end{equation*}
    where $X$ is observed. Here, $f$ is locally, but not globally, identifiable.
    \begin{proof}[Proof]
        If two admissible functions $f_1$ and $f_2$ give the same observed dynamics, then
        \begin{equation*}
            f_1(x)^2 = f_2(x)^2.
        \end{equation*}
        Here, both the alternatives $f_1=f_2$ and $f_1=-f_2$ may hold, implying that at least two potential functions will be consistent with any given observed dynamics. This implies local functional identifiability.
    \end{proof}

    We note that $f$ is globally identifiable at any point where $f_1(x^*)=f_2(x^*)=0$. Furthermore, even if $f$ is assumed to be continuous, whichever of the relations $f_1=f_2$ and $f_1=-f_2$ holds may switch at points where $f_1(x^*)=f_2(x^*)=0$, permitting the generation of a potentially large equivalence class of functions.

    \item \textbf{Function-label permutation ambiguity.} Consider the two-state model
    \begin{align*}
        \dot{X}_1 &= f(X_1)+g(X_1),\\
        \dot{X}_2 &= f(X_1)g(X_1),
    \end{align*}
    where both $X_1$ and $X_2$ are observed. Here, $f$ and $g$ are locally, but not globally, identifiable.
    \begin{proof}[Proof]
        Let $(f_1,g_1)$ and $(f_2,g_2)$ be two admissible pairs producing the same observed dynamics. Since both state variables are observed, the model equations give
        \begin{align*}
            f_2(x)+g_2(x)=f_1(x)+g_1(x),
            \qquad
            f_2(x)g_2(x)=f_1(x)g_1(x),
        \end{align*}
        on the observed range of $X_1$. From the first relation, $g_2(x)=f_1(x)+g_1(x)-f_2(x)$. Substituting this into the second relation gives
        \begin{equation*}
            \bigl(f_2(x)-f_1(x)\bigr)\bigl(f_2(x)-g_1(x)\bigr)=0.
        \end{equation*}
        Here, both the alternatives $(f_1,g_1)=(f_2,g_2)$ and $(f_1,g_1)=(g_2,f_2)$ may hold, implying that at least two possible pairs of functions will be consistent with any given observed dynamics. This implies local functional identifiability.
    \end{proof}

    We note that $(f,g)$ is globally identifiable at any point where $f_1(x^*)=g_1(x^*)$. Furthermore, even if $f$ and $g$ are assumed to be continuous, whichever of the relations $(f_1,g_1)=(f_2,g_2)$ and $(f_1,g_1)=(g_2,f_2)$ holds may switch at points where $f_1(x^*)=g_1(x^*)$, permitting the generation of a potentially large equivalence class of function pairs.

    \item \textbf{Argument-reflection ambiguity.} Consider
    \begin{align*}
        \dot{X}_1 &= f(X_1)+f(-X_1),\\
        \dot{X}_2 &= f(X_1)f(-X_1),
    \end{align*}
    where both states are observed. Here, $f$ is locally, but not globally, identifiable.
    \begin{proof}[Proof]
        Let $f_1$ and $f_2$ give the same observed dynamics. Then
        \begin{align*}
            f_2(x)+f_2(-x)=f_1(x)+f_1(-x),
            \qquad
            f_2(x)f_2(-x)=f_1(x)f_1(-x),
        \end{align*}
        for the observed values $x$ of $X_1$. From the first relation, $f_2(-x)=f_1(x)+f_1(-x)-f_2(x)$. Substituting this into the second relation gives
        \begin{equation*}
            \bigl(f_2(x)-f_1(x)\bigr)\bigl(f_2(x)-f_1(-x)\bigr)=0.
        \end{equation*}
        Here, both the alternatives $f_1=f_2$ and $f_2(x)=f_1(-x)$ may hold, implying that at least two possible functions will be consistent with any given observed dynamics. This implies local functional identifiability.
    \end{proof}

    We note that $f$ is globally identifiable at any point where $f_1(x^*)=f_1(-x^*)$. Furthermore, even if $f$ is assumed to be continuous, whichever of the relations $f_1(x^*)=f_2(x^*)$ and $f_2(x^*)=f_1(-x^*)$ holds may switch at points where $f_1(x^*)=f_1(-x^*)$, permitting the generation of a potentially large equivalence class of functions.

    \item \textbf{Input-coordinate ambiguity.} Consider
    \begin{align*}
        \dot{X}_1 &= f(X_1,X_2)+f(X_2,X_1),\\
        \dot{X}_2 &= f(X_1,X_2)f(X_2,X_1),
    \end{align*}
    where both states are observed. Here, $f$ is locally, but not globally, identifiable.
    \begin{proof}[Proof]
        Let $f_1$ and $f_2$ give the same observed dynamics. Since both states are observed,
        \begin{align*}
            f_2(x,y)+f_2(y,x)=f_1(x,y)+f_1(y,x),
            \qquad
            f_2(x,y)f_2(y,x)=f_1(x,y)f_1(y,x)
        \end{align*}
        on the observed pairs $(x,y)=(X_1,X_2)$. From the first relation, $f_2(y,x)=f_1(x,y)+f_1(y,x)-f_2(x,y)$. Substituting this into the second relation gives
        \begin{equation*}
            \bigl(f_2(x,y)-f_1(x,y)\bigr)\bigl(f_2(x,y)-f_1(y,x)\bigr)=0.
        \end{equation*}
        Here, both the alternatives $f_1=f_2$ and $f_2(x,y)=f_1(y,x)$ may hold, implying that at least two possible functions will be consistent with any given observed dynamics. This implies local functional identifiability.
    \end{proof}

    We note that $f$ is globally identifiable at any point where $f_1(x^*,y^*)=f_1(y^*,x^*)$. Furthermore, even if $f$ is assumed to be continuous, whichever of the relations $f_1(x,y)=f_2(x,y)$ and $f_2(x,y)=f_1(y,x)$ holds may switch at points where $f_1(x^*,y^*)=f_1(y^*,x^*)$, permitting the generation of a potentially large equivalence class of functions.
\end{enumerate}

% ------------------------------------------------------------------------------------------------------------------------------------ %
% ------------------------------------------------- Appendix D: chemical reaction network model derivation ------------------------------------------------- %
% ------------------------------------------------------------------------------------------------------------------------------------ %

\section{Chemical reaction network model derivation} \label{appendix:chemical reaction network_derivation}

A chemical reaction network model consists of a set of species and a set of reactions. Here we consider a model with species $X_1$ and $X_2$ and the following reaction events:
\begin{align*}
\ce{ $X_1$ &->C[$1$] $X_1 + X_2$, } \\
\ce{ $X_2$ &->C[$g(X_2)$] $X_1$,} \\
\ce{ $X_1$ &->C[$d$] \text{\o},}
\end{align*}
where $d>0$ is a parameter and $g$ is a function of the concentration of $X_2$. Application of the law of mass action, gives the following ODE model:
\begin{align*}
    \dot{X}_1 = g(X_2)X_2 - dX_1, \\
    \dot{X}_2 = X_1 - g(X_2)X_2.
\end{align*}
Setting $f(X_2) = g(X_2)X_2$ generates the model in Equations~\eqref{eq:chemical reaction network_dXdt}--\eqref{eq:chemical reaction network_dYdt}.

% ------------------------------------------------------------------------------------------------------------------------------------ %
% ------------------------------------------------- Appendix E: Numerical demonstration ------------------------------------------------- %
% ------------------------------------------------------------------------------------------------------------------------------------ %

%\section{Numerical demonstration} \label{sec:appendix_numerics}
%I will add this for as many models as we can be bothered with and that we want. In principle, we can do for all. We probably should do for a few at least. We might want to expand this a bit between preprint and submission.

% [Numerical demonstration of the results in Sections~\ref{sec:diffalg}--\ref{sec:applications} to be added.]

% ------------------------------------------------------------------------------------------------------------------------------------ %

\end{document}
